% !TEX root = nc_cubic_threefold_moduli.tex
\documentclass[11pt,a4paper]{ltjsarticle}
\usepackage{amsmath,amssymb,amsthm,mathtools,bm,mathrsfs}
\usepackage{luatexja-fontspec}
\usepackage{enumitem,booktabs,array}
\usepackage{xcolor}
\usepackage[colorlinks=true,linkcolor=blue!45!black,citecolor=blue!45!black,urlcolor=blue!45!black]{hyperref}
\IfFileExists{/usr/share/fonts/opentype/noto/NotoSerifCJK-Regular.ttc}{%
  \setmainjfont[Path=/usr/share/fonts/opentype/noto/]{NotoSerifCJK-Regular.ttc}%
}{\setmainjfont{HaranoAjiMincho-Regular}}
\setlength{\textwidth}{15.6cm}
\setlength{\oddsidemargin}{0.25cm}
\setlength{\evensidemargin}{0.25cm}
\setlength{\textheight}{22.6cm}
\setlength{\topmargin}{-0.2cm}
\setlength{\parindent}{1em}
\setlength{\parskip}{1pt}
\numberwithin{equation}{section}
\newtheorem{theorem}{定理}[section]
\newtheorem{lemma}[theorem]{補題}
\newtheorem{proposition}[theorem]{命題}
\newtheorem{corollary}[theorem]{系}
\newtheorem{remark}[theorem]{注意}
\newtheorem{definition}[theorem]{定義}
\renewcommand{\proofname}{証明}
\newcommand{\C}{\mathbb C}
\providecommand{\CC}{\C}
\newcommand{\PP}{\mathbb P}
\newcommand{\A}{\mathbb A}
\newcommand{\OO}{\mathcal O}
\newcommand{\cB}{\mathcal B}
\newcommand{\cA}{\mathcal A}
\newcommand{\cI}{\mathcal I}
\newcommand{\cQ}{\mathcal Q}
\newcommand{\cZ}{\mathcal Z}
\newcommand{\Sym}{\operatorname{Sym}}
\newcommand{\gr}{\operatorname{gr}}
\newcommand{\qgr}{\operatorname{qgr}}
\newcommand{\gldim}{\operatorname{gldim}}
\newcommand{\rank}{\operatorname{rank}}
\newcommand{\Span}{\operatorname{span}}
\newcommand{\Spec}{\operatorname{Spec}}
\newcommand{\Aut}{\operatorname{Aut}}
\newcommand{\Hom}{\operatorname{Hom}}
\newcommand{\Ext}{\operatorname{Ext}}
\newcommand{\Proj}{\operatorname{Proj}}
\newcommand{\Frac}{\operatorname{Frac}}
\newcommand{\trdeg}{\operatorname{trdeg}}
\newcommand{\GL}{\mathrm{GL}}
\newcommand{\PGL}{\mathrm{PGL}}
\newcommand{\Lie}{\operatorname{Lie}}
\newcommand{\Z}{\operatorname{Z}}
\newcommand{\Tens}{\operatorname{T}}
\newcommand{\CY}{\mathrm{CY}}
\newcommand{\AS}{\mathrm{AS}}
\newcommand{\nc}{\mathrm{nc}}
\newcommand{\Symm}{\operatorname{sym}}
\title{非可換三次三様体の17次元モジュライ族\\中心三次対と変形幾何}
\author{公表準備用研究原稿（著者情報未記入）}
\date{2026年11月13日・nc cubic moduli研究草稿 v5}
\begin{document}
\maketitle
\begin{abstract}
相対Jacobian型の中心三次対の41次元既約代数族を構成し、
18次元の放物型群作用に横断する17係数正規形を与える。
31次整数行列の非零行列式によって横断性を証明し、
二次包絡による同型類判定から実際の非可換三次三様体の
次数付き同型類が17次元分だけ変動することを示す。
さらに分離型の余次元15、可換cubicへの忘却核7次元の幾何学的分解、
混合中心元と相対Picardの障害を扱う。
また、正則正規$d$次超曲面を持つKoszul AS正則$\mathbb Z$-代数の
最小周期が$\gcd(N,d)$を割ることを示し、5次元・三次の場合の1周期性を導く。
この結論はambientの周期性であり、商イデアルを保つ1周期構造は主張しない。
一般次元のambient復元とその証明はnc-cubicの論文Aにも収録する。
滑らかな可換三次点の全Artin形式局所環とOre局所化による
中心三次対・構造代数層・斉次環の比較を、独立の証明依存を明示して収録する。
また低次数中心・偶数Veronese・構造対象付き導来復元と、
具体的な剛体ambient上の四次元非同型族を条件を明記して扱う。
構成した族が全nc cubic moduliの既約成分をなすことは未証明である。
\end{abstract}
\tableofcontents

\section{Introduction：主張と先行研究}\label{sec:intro}
非可換射影空間を次数付き AS 正則代数の $\qgr$ として扱う立場は
Artin--Schelter \cite{AS} と Artin--Zhang \cite{AZ} の枠組みに属する．
可換な cubic threefold $X\subset\PP^4$ は５変数多項式環の三次正則元による
商として記述される．その非可換類似として，次の組を考える：
\begin{equation}\label{eq:data}
 (B,[G]),\qquad B=\Tens(V)/(R_2),\quad\dim V=5,\quad
 G\in\Z(B)_3\setminus\{0\},\quad A=B/(G).
\end{equation}
ここで $B$ は連結次数付き，noetherian，Koszul，AS 正則次元５であり，
$[G]$ は非零スカラー倍を同一視した中心正則元の直線である．
本稿はこの型の\emph{一つの具体的な既約族の一般点}を扱い，
$A$ から $B$ を内在的に復元する．任意の非可換三次三様体が
\eqref{eq:data} の形に書けるとは仮定も主張もしない．

\paragraph{半古典極限と既知の分類.}
Loray--Pereira--Touzet \cite[Theorem~9.1]{LPT} と
Mayanskiy \cite[Theorem~1]{Mayanskiy} は，固定した滑らかな cubic threefold の
射影 Poisson 空間が Pl\"ucker 埋め込みされた $\operatorname{Gr}(2,5)$ であることを示した．
その結果は「固定した可換 cubic の方向が７次元」であることを示唆するが，
17次元の\emph{有限係数の非可換代数族}を直接与えるわけではない．
Pym \cite[Proposition~2.9, Theorem~2.10]{Pym} は，射影空間の
Calabi--Yau 型の量子変形と二次 Poisson 構造の成分の関係を証明し，
非可換 $\PP^3$ の場合に具体的な分類を与えた．本稿は
「周囲の量子構造に中心三次方程式を指定する」という別の incidence を対象にする．

\paragraph{代数構成の入力.}
三変数 Sklyanin 代数と楕円曲線の対応，次数３の中心元および
twisted homogeneous coordinate ring については
Artin--Tate--Van den Bergh \cite{ATV} が基本的な文献である．
Mori--Smith \cite[Theorem~4.2(1)]{MS} によって，
滑らかな平面三次式を用いた Jacobian 型二次関係式は
３次元 Calabi--Yau／AS 正則代数を与える．
Berger--Taillefer \cite{BT} は低次 potential の PBW 変形を扱い，
Etingof--Ginzburg \cite[Proposition~7.3.1(ii), Theorem~3.3.2]{EG} は
それを Calabi--Yau 性と中心元の存在に結び付けた．
同論文付録の Rains による $\widetilde E_6$ 型中心元の有限公式も重要な入力である
\cite[\S9.1]{EG}．以下ではこの文献群を二つの中心一次変数に相対化する．

\paragraph{復元の位置付け.}
Artin--Zhang \cite{AZ} が扱う偏極付き三つ組の section algebra から
次数付き代数を回復する一般原理と，
「三次関係は二次関係を変えない」という初等的な観察を組み合わせると，
ambient の quadratic envelope が得られる．一方，構造対象だけを残して
次数シフトを忘れる場合には，Yekutieli--Zhang \cite{YZ} の Serre 双対に基づく
偶数 Veronese の復元が使える．この後半には明示的な追加仮定が必要である．

本稿で扱う新規性の候補は，\emph{全混合係数を許す相対中心三次族，
17係数の横断的正規形，およびそれと二次包絡の一体化}にある．
個別の結論が先行研究で未発表であることまで確認したわけではない．
特に AS 正則性の一般論，PBW 定理，中心元の原公式，section algebra の復元自体は
既知の結果として明記する．
分離した一楕円ブロックは中心対の有理商で余次元15を占めるため、
その明示的可換復元を一般17次元族へ無条件に延長することはできない
（命題\ref{prop:split-codimension}）。

\paragraph{一般のAS正則$\mathbb Z$-代数との関係.}
正則正規な三次超曲面を持つ標準5次元Koszul AS正則$\mathbb Z$-代数は、
AS双対性の5周期と正規族の3周期により1周期的である
（系\ref{cor:moduli-cubic-ambient-periodic}）。
これは通常の次数付きambientを仮定しない場合の結論であるが、
1周期同型が超曲面イデアルを保つことは別問題である。
既知のKoszul--Frobenius双対性\cite{PBondalPolishchukHelices}と
正規族\cite{PChanNymanNormalFamily}に基づき、初出の網羅的調査は未完了である。

\paragraph{主定理.}
第\ref{sec:construction}節で中心三次対の41次元族、
第\ref{sec:normal}節で17係数の横断形、
第\ref{sec:split-codimension}節で分離型と忘却核、
第\ref{sec:mixed-quantum-center}節で混合中心・相対Picardの障害、
第\ref{sec:moduli-iso}節で同型類の判定を扱う。

\section{相対 Jacobian 型の中心三次族}\label{sec:construction}
本節では基礎体を $k=\C$ とする．$U=ku\oplus kv$ と
$W=kx\oplus ky\oplus kz$ を置き，すべての生成元に次数１を与える．
$F\in\Sym^3(W\oplus U)$ を
\begin{equation}\label{eq:Ffull}
 F=f_3(x,y,z)+Q_2(x,y,z;u,v)+L_1(x,y,z;u,v)+g(u,v)
\end{equation}
と書く．$f_3$ は $x,y,z$ に関して次数３，$Q_2$ は次数２，
$L_1$ は次数１，$g\in\Sym^3U$ であり，それぞれ全次数３で斉次とする．

\begin{definition}\label{def:B}
$\Symm$ を可換単項式の完全対称化とする．$B_F$ を
\begin{equation}\label{eq:B}
 B_F=\frac{k\langle x,y,z,u,v\rangle}{
 \bigl([u,-],[v,-],\ [x,y]-\Symm(F_z),\
 [y,z]-\Symm(F_x),\ [z,x]-\Symm(F_y)\bigr)}
\end{equation}
で定義する．$g(u,v)$ は右辺の微分に現れない．
\end{definition}

\begin{lemma}[平面三次式の相対 PBW 性]\label{lem:pbw}
$E=V(f_3)\subset\PP(W^*)$ が滑らかなら，$B_F$ は
$k[u,v]$ 上自由であり，$x,y,z$-次数による濾過について
\begin{align}\label{eq:gr}
 \gr_{x,y,z}B_F&\simeq S_{f_3}\otimes_k k[u,v],\\
 S_{f_3}&:=k\langle x,y,z\rangle/\bigl([x,y]-\Symm((f_3)_z),\ 
\notag\\[-2pt]
 &\hspace{8em}[y,z]-\Symm((f_3)_x),\ [z,x]-\Symm((f_3)_y)\bigr).
\end{align}
さらに $S_{f_3}$ は noetherian な３次元 Koszul AS 正則 Calabi--Yau 整域で，
$B_F$ は noetherian な５次元 Koszul AS 正則整域となる．とくに
\begin{equation}\label{eq:HilbertB}
 H_{B_F}(t)=(1-t)^{-5}.
\end{equation}
\end{lemma}
\begin{proof}
Mori--Smith \cite[Theorem~4.2(1)]{MS} により，滑らかな $f_3$ から得た
$S_{f_3}$ は３次元 Calabi--Yau であり，したがって上記の
AS 正則性と Koszul 性を持つ．$u,v$ を中心的な係数とみなし，
\eqref{eq:B} の三つの関係を
$S_{f_3}$ の potential に対する次数２以下の cyclic potential の摂動と解釈する．
各 $(u,v)=(s,t)\in k^2$ への特殊化について，
Berger--Taillefer および Etingof--Ginzburg の PBW 定理
\cite[Proposition~7.3.1(ii)]{EG} により，$x,y,z$-濾過の随伴次数付き代数は
$S_{f_3}$ に等しい．

$S_{f_3}$ の斉次な語の基底を選ぶ．元の関係式を用い，語を最高 $x,y,z$-次数から
下げる帰納法により，選んだ語の持上げは $k[u,v]$ 上 $B_F$ を生成する．
この持上げの有限 $k[u,v]$-線形関係があれば，各 $(s,t)$ への特殊化で
PBW 基底間の関係を与えるので，その全係数は $k[u,v]$ の全閉点で零となり，
したがって零である．これで自由性および \eqref{eq:gr} が得られる．
特に $u,v$ は $B_F$ 上の正則中心列である．

$B_F/(u,v)\simeq S_{f_3}$ なので，中心正則元に関する
AS 正則性の上昇を二度適用すれば次元５の AS 正則性が得られる
\cite[Propositions~2.1, 2.3]{CKS}．noetherian 性は中心正則拡大の
標準定理，Koszul 性は中心一次元による change-of-rings，整域性は
\eqref{eq:gr} の整域性に従う．$u,v$ を加えた次数計数から \eqref{eq:HilbertB} を得る．
\end{proof}

\begin{lemma}[ambient の Calabi--Yau 性]\label{lem:CY}
補題\ref{lem:pbw} の仮定の下で $B_F$ の中山自己同型は恒等であり，
$B_F$ は５次元 Calabi--Yau 代数である．
\end{lemma}
\begin{proof}
AS 正則性から $B_F$ は skew Calabi--Yau であり，中山自己同型を $\nu$ とする．
中山双加群の左右の中心作用は一致するため，$\nu$ は中心 $k[u,v]$ を
各点固定する．すなわち $\nu(u)=u,\nu(v)=v$ である．
したがって各 $(s,t)\in k^2$ について $\nu$ は
$B_F/(u-s,v-t)$ に自己同型を誘導する．
Zhou--Shen--Lu の正規元による中山自己同型の特殊化公式
\cite[Propositions~2.6, 3.6]{ZSL} により，この自己同型は
特殊化代数の中山自己同型である．
一方，特殊化代数は低次 potential の摂動であり，
Etingof--Ginzburg の定理\cite[Theorem~3.3.2(i)]{EG} によって
３次元 Calabi--Yau である．特殊化代数は随伴次数付き整域
$S_{f_3}$ を持つので，単元は $k^\times$ のみである．
従って特殊化された $\nu$ は恒等である．
$B_F$ は $k[u,v]$ 上自由だから，
$\nu(x)-x,\nu(y)-y,\nu(z)-z$ が全 $(s,t)$ で零になることから
それらは零である．よって $\nu=\mathrm{id}$．
\end{proof}

\begin{lemma}[一般点での三次中心元]\label{lem:center}
$E=V(f_3)$ が滑らかな開集合のある非空開部分で
\begin{equation}\label{eq:centS}
 \Z(S_{f_3})_1=\Z(S_{f_3})_2=0,\qquad
 \dim_k\Z(S_{f_3})_3=1
\end{equation}
が成立する．その開部分をさらに縮小し，三次中心の非零元
$\Omega_{f_3}\in\Z(S_{f_3})_3$ を代数的に選ぶ．
すると\eqref{eq:B} の $Q_2,L_1$ の全係数に関して代数的に変わる
\begin{equation}\label{eq:G}
 \mathcal G_F=\Omega_{f_3}+C_2(F)+C_1(F)\in\Z(B_F)_3
\end{equation}
が存在する．$C_2$ は $x,y,z$-次数２で $u,v$ に線形，$C_1$ は
$x,y,z$-次数１で $u,v$ に二次であり，$x,y,z$-次数０の部分を零とする
規格化によって一意である．特に
\begin{equation}\label{eq:Gg}
 G_F=\mathcal G_F+g(u,v)
\end{equation}
も中心元であり，$B_F$ 上正則である．
\end{lemma}
\begin{proof}
具体的に $f_3=(x^3+y^3+z^3)/3$ では，
$S_{f_3}=k\langle x,y,z\rangle/([x,y]-z^2,[y,z]-x^2,[z,x]-y^2)$ であり，
自由語の次数１から４の関係行列を有理数上で消去すると，
次数１，２，３の中心の次元は順に $0,0,1$ になる．
この有限計算は付属 \nolinkurl{computations/independent/verify_generic17_center.py} で再現できる．
三変数 Sklyanin 代数の楕円曲線と translation の表示
\cite{ATV,EG} において，$\tau,\tau^2\ne1$ となる一般の点では
低次数中心が零であり，中心三次元の空間は一次元である．
この条件を中心性の次数４以下の有限行列の rank として書けば
非空 Zariski 開条件になり，その上では \eqref{eq:centS} の
核は階数１の代数的線束をなす．非零小行列式を反転した
chart 上で $\Omega_{f_3}$ を選べる．

$u,v$ を数値へ特殊化した三変数の非斉次 potential を考える．
\cite[Theorem~3.3.2(ii)]{EG} は $x,y,z$-濾過次数３以下に
非スカラー中心元が存在することを保証する．その最高次数の記号は
\eqref{eq:centS} により $\Omega_{f_3}$ の非零定数倍であり，
係数を規格化できる．補正項の $x,y,z$-次数２および１の係数は
それぞれ６個および３個の未知数を持ち，中心性は線形方程式
\begin{equation}\label{eq:centerlinear}
 M(\mathbf a,\mathbf b)c=d(\mathbf a,\mathbf b)
\end{equation}
で表される．ここで $\mathbf a$ は非斉次二次 potential の係数，
$\mathbf b$ は一次 potential の係数である．
零定数部分で \eqref{eq:centerlinear} の同次解が非零なら，
$x,y,z$-濾過次数２以下の非スカラー中心元を与え，
最高記号を取ると \eqref{eq:centS} に矛盾する．
よって $M$ は全特殊化で階数９を持つ．各特殊化に解が存在するので，
一定 rank のベクトル束写像の余核中で $d$ の像は全閉点で零となる．
基底を被約な多項式係数環としたからその像は零であり，
唯一の代数的な解が存在する．

$\mathbf a$ に重み１，$\mathbf b$ に重み２を与える
$\mathbb G_m$-作用と規格化の一意性から，$C_2$ は
$\mathbf a$ に一次，$C_1$ は $\mathbf a$ に二次・$\mathbf b$ に一次である．
$\mathbf a$ を $u,v$ の一次式，$\mathbf b$ を二次式に置けば
\eqref{eq:G} の斉次中心元を得る．
$\Omega_{f_3}$ は \eqref{eq:gr} の整域における非零元なので，
$G_F$ は $B_F$ の非零正則元である．
\end{proof}

\begin{theorem}[41次元の中心三次対の代数的族]\label{thm:incidence}
$U=\langle u,v\rangle$ を固定し，$F$ の35係数を動かす．
補題\ref{lem:center} の非空 Zariski 開条件を課したパラメータ空間
$\mathcal T_U$ 上，$F\mapsto(B_F,[G_F])$ は代数的な中心三次対の族となる．
商
\begin{equation}\label{eq:A}
 A_F=B_F/(G_F)
\end{equation}
は各次数で基底上局所自由であり，
\begin{equation}\label{eq:HseriesA}
 H_{A_F}(t)=\frac{1-t^3}{(1-t)^5},\qquad
 (\dim A_1,\dim A_2,\dim A_3)=(5,15,34).
\end{equation}
$U$ を $\operatorname{Gr}(2,5)$ 上で動かすと，この族のパラメータ空間は
$6+35=41$ 次元の既約な有限型空間を与える．その ambient 二次代数の
係数空間は $6+31=37$ 次元であり，一般の固定 ambient の上で
中心三次直線の選択に４次元の自由度がある．
\end{theorem}
\begin{proof}
一般の斉次三次式の係数は $\binom{7}{3}=35$ 個であり，
$F$ に含まれる $g\in\Sym^3U$ の係数は４個である．
$B_F$ の関係式は $F$ の $x,y,z$ 偏微分から決まり，$g$ には依存しない．
補題\ref{lem:center} の $\mathcal G_F$ は最高 $x,y,z$-次数成分によって規格化され，
$G_F=\mathcal G_F+g$ であるから，$g$ の異なる４係数は
同じ ambient 内の異なる中心三次直線を与える．
補題\ref{lem:pbw} および補題\ref{lem:center} により，$B_F$ は整域，$G_F$ は
中心正則である．$B_F$ の任意の次数に対して中心正則元による完全列
\[
 0\longrightarrow B_F(-3)\xrightarrow{\cdot G_F}
 B_F\longrightarrow A_F\longrightarrow0
\]
が成立し，\eqref{eq:HseriesA} が従う．
相対的には，固定次数で有限生成自由加群間の表示行列の rank が
全閉点で一定なので，次数成分は局所自由となる．
$U$ の変化は Grassmannian 上の局所 chart ごとに同じ構成で与えられ，
全体の既約次元は $\dim\operatorname{Gr}(2,5)+35=41$ である．
\end{proof}

\begin{remark}[正則性の二つの意味]\label{rem:regularity}
本定理は ambient $B_F$ の５次元 AS 正則性，$G_F$ の正則性，
$A_F$ の相対平坦性を示す．これだけで任意の閉点について
$\gldim(\qgr A_F)=3$ が従うわけではない．
量子的な楕円境界が滑らかな固定 ambient について，十分一般の
中心三次係数で tails が滑らかになるという議論には，
中心局所化と非可換 generic smoothness の結果
\cite{RSVdB,Rains} が必要である．
本稿の17次元性および復元定理は，この追加の全閉点での
Zariski 開性を仮定しない．
\end{remark}

\section{generic な17係数正規形と真の変形次元}\label{sec:normal}
ここで重要なのは，41係数から形式的に24を引くことだけではない．
実際の座標変換の一般軌道が24次元であり，17個の係数だけを持つ族が
その有理商を支配することを示す．

\subsection{中心一次平面と群作用}
\begin{lemma}[一次中心による平面の復元]\label{lem:U}
補題\ref{lem:center} の開集合では
\begin{equation}\label{eq:Z1}
 \Z(B_F)_1=U=ku\oplus kv.
\end{equation}
したがって，次数付き代数 $B_F$ の自己同型は $U$ を保存する．
\end{lemma}
\begin{proof}
$u,v$ は中心である．$w\in (B_F)_1$ が中心なら，
$x,y,z$-次数に関する最高記号は $S_{f_3}[u,v]$ の中心である．
その $x,y,z$-次数が１なら，$\Z(S_{f_3})_1=0$ に矛盾する．
よって $w\in U$ である．
\end{proof}

$P\subset\PGL(W\oplus U)$ を平面 $U$ の保存群とする．
標識された基底 $(x,y,z;u,v)$ に関して，$P$ は
\[
 \left(\begin{array}{cc}M&C\\0&N\end{array}\right),\qquad
 M\in\GL_3,\quad N\in\GL_2,\quad C\in\operatorname{Mat}_{3\times2}
\]
の同時スカラー倍を同一視した群である．したがって
\begin{equation}\label{eq:dimP}
 \dim P=9+4+6-1=18,
 \qquad\dim\PGL_5=24.
\end{equation}
Jacobian 括弧は変数変換に対し $\det M$ の重みを持つ．
したがって ambient 二次関係式の係数 $F$ の
$\Sym^3U$ 以外の31成分への Lie 作用は
\begin{equation}\label{eq:Lie}
 \delta_{M,N,C}F=(\operatorname{tr}M)F
 -\sum_{i=1}^{3}\bigl((Mx)_i+(Cs)_i\bigr)\partial_{x_i}F
 -\sum_{j=1}^{2}(Ns)_j\partial_{s_j}F,
 \qquad s=(u,v)^t,
\end{equation}
を $\Sym^3U$ で割った像である．
\emph{中心元の規格化は変数変換で binary cubic の補正を受け得る}ため，
以下の横断性は最初に ambient の31方向で確かめ，
最後に独立な中心三次の４方向を加える．
この順序なら中心元の gauge に関する余分な同変性は不要である．

\subsection{17係数の表示}
\begin{theorem}[17係数の横断的代数的族]\label{thm:17}
$M_1,M_2,M_3\in k[u,v]_2$, $g\in k[u,v]_3$ として
\begin{equation}\label{eq:F17}
\boxed{\begin{aligned}
 F_{17}={}&a(x^3+y^3+z^3)+bxyz\\
 &+ux^2+\alpha vy^2+\beta(u+v)z^2\\
 &+M_1(u,v)x+M_2(u,v)y+M_3(u,v)z+g(u,v).
\end{aligned}}
\end{equation}
この族の係数は
\begin{equation}\label{eq:count17}
 (a,b,\alpha,\beta)\quad\text{の４個},\qquad
 (M_1,M_2,M_3)\quad\text{の９個},\qquad
 g\quad\text{の４個}
\end{equation}
であり，合計17個である．
非空 Zariski 開集合 $T\subset\A^{17}$ 上で
$F=F_{17}$ を \eqref{eq:B}, \eqref{eq:Gg} に代入すると
$A_F=B_F/(G_F)$ は次数ごとに平坦な代数的族となる．
その一般点の次数付き代数同型類は17次元分だけ動く．
より正確には，この族から第\ref{sec:construction}節の
中心三次対の Rosenlicht 有理商への有理写像は支配的かつ一般に有限で，
\begin{equation}\label{eq:17mod}
 \trdeg_k k(\cI)^{\PGL_5}=41-24=17.
\end{equation}
ここで $\cI$ は構成した中心三次対の41次元既約族の像の閉包とする．
\end{theorem}
\begin{proof}
まず35係数空間での一般の cubic $F$ の $f_3$ が滑らかな
平面三次式を定めるとする．古典的 Hesse 正規形
\cite[\S2, Lemma~2.1]{AD} により，
$\GL_3$ の座標変換で $f_3=a(x^3+y^3+z^3)+bxyz$ と書ける．
$b\ne0$ とし，$x,y,z$ を $u,v$ の一次式だけ平行移動すると，
三つの交差二次項 $xy,yz,zx$ に付く二元一次式６係数を消せる．
残る３つの二元一次式が $\PP(U)$ 上に互いに異なる零点を持つ一般点では，
$\PGL(U)$ でそれらを $u,v,u+v$ に変換し，
相対スカラーで第一の係数を１に規格化すれば \eqref{eq:F17} になる．
ただし Hesse 変換の選択は有限多価であり，この議論だけで
有理商の一般次数が１であるとはいわない．

さらに代数的な横断性を有理数体で独立に検算する．
\eqref{eq:F17} の次の点を取る：
\begin{align}\label{eq:F0}
 F_0={}&3(x^3+y^3+z^3)+5xyz+ux^2+2vy^2+3(u+v)z^2\notag\\
 &+(u^2+uv+2v^2)x+(2u^2+3uv+v^2)y\\
 &+(u^2+4uv+5v^2)z+u^3+u^2v+2uv^2+3v^3.\notag
\end{align}
$\mathfrak{gl}_3\oplus\mathfrak{gl}_2\oplus\operatorname{Mat}_{3\times2}$
の19個の行列係数に対する微分 \eqref{eq:Lie} の像を取り，
これと \eqref{eq:F17} の方向のうち $g$ 以外の13方向とを
$\Sym^3(W\oplus U)/\Sym^3U$ の31個の単項式に展開する．
同時スカラー１方向を除いた $31\times31$ の行列について
\begin{equation}\label{eq:det}
 \boxed{\det M_{31}=-5538243375000\ne0}
\end{equation}
となる．行と列の取り方を付録\ref{app:rank}に固定した．
この計算は有理数体上の厳密計算であり，
付属 \nolinkurl{computations/independent/verify_generic17_slice.py} で再現できる．
したがって ambient 二次代数の31方向において，
$P$ の一般軌道は18次元で，13係数正規形は横断的である．

残りの $\Sym^3U$ の４係数は，固定 ambient の上の
$[\mathcal G_F+g]$ の４次元のアフィン chart である．
よって全中心三次対の固定 $U$ での一般軌道も18次元であり，
17係数正規形は一般軌道に横断する．
$U$ が \eqref{eq:Z1} によって intrinsic なので，
一般の同型がこれ以外の６次元の pencil 方向を
同一視してしまうことはない．
$U$ を動かすと $\cI$ の次元は41，$\PGL_5$ の一般軌道の次元は
$6+18=24$ である．Rosenlicht の有理商の定理 \cite{Rosenlicht}
と次元公式により \eqref{eq:17mod} を得る．
\end{proof}

\begin{remark}[定義域と17次元の意味]\label{rem:meaning}
この次元は $\PGL_5$ による同型軌道についての値であり，
実際に $A$ の次数付き同型類に対応することは
次節の復元定理\ref{thm:moduli-iso}で保証する．
$\PGL_5$ の一般 stabilizer が有限であることは示したが，
自明であることや有限次数が１であることは示していない．
また $T$ の各点で $\qgr A$ が滑らかであるとは主張しない．
\end{remark}

\section{分離型の余次元と一般混合項の不可避性}
\label{sec:split-codimension}
一楕円Sklyaninブロックと独立な中心三次式からなる分離型は、
rotation付き可換復元の明示的な成功例である。
しかし、それが本稿の17次元一般族を支配するわけではない。
その理由を中心一次平面を固定した係数空間で正確に述べる。

\begin{proposition}[分離型は余次元15]
\label{prop:split-codimension}
$W=\langle x,y,z\rangle$、$U=\langle u,v\rangle$、
$\mathcal T=\Sym^3(W\oplus U)$とする（$\dim\mathcal T=35$）。
$W$上の滑らかな三次形式$f$、$U$上の一般の三次形式$g$、
および$L\in\Hom(U,W)$からなる
\[
 F_{f,g,L}(x,u)=f(x+L(u))+g(u)
 \tag{S}\label{eq:split-family}
\]
の像$\mathcal S\subset\mathcal T$は既約な20次元構成可能集合
の閉包を持ち、$\mathcal T$内で余次元15である。
さらに、中心平面$U$を保つ本稿の放物型座標変換群
$\mathcal P$の一般軌道が18次元であるため、
$\mathcal S$の有理商は次元2である。
従って、一般17次元の中心三次対の有理商において、
分離型の像は余次元15の真の特殊部分をなす。
\end{proposition}
\begin{proof}
$f$の係数は10、$g$は4、$L$は6で、全20変数である。
$L=0$における微分は$f$と$g$の係数方向のほか、
\[
 u f_x,\ u f_y,\ u f_z,\quad
 v f_x,\ v f_y,\ v f_z
 \tag{T}\label{eq:split-tangent}
\]
を与える。$f$が滑らかなら$f_x,f_y,f_z$は定数係数で一次独立である。
実際、定数線形関係があれば$f$はある座標に依存しない錐となり、
射影平面三次曲線は特異になる。
従って\eqref{eq:split-tangent}の六方向は独立で、
$\Sym^3W$と$\Sym^3U$の14方向とも次数で分離される。
ゆえに微分の階数は20であり、$\dim\overline{\mathcal S}=20$。

群$\mathcal P$は、$W$内と$U$内の線形座標変換、
$W$への$U$のshearを含み、$F$への作用は本稿の
Poisson中心三次対の規約（式\eqref{eq:Lie}）による。
この群は$\mathcal S$を保つ。
一般の滑らかな$f$と重根のない二元三次式$g$は、
それぞれ射影的な安定化群が有限である。
$L=0$で$f$を固定するshearの無限小方向は
\eqref{eq:split-tangent}の独立性から消える。
従って$\mathcal P$の有効な一般安定化群は有限であり、
一般軌道の次元は18となる。
ゆえに有理商の次元は$20-18=2$。
全係数空間の一般有理商は第\ref{sec:normal}節により
$35-18=17$次元なので、その差は15である。
\end{proof}

\begin{corollary}[本質的な混合項の一次近似]
\label{cor:split-normal-directions}
分離点$F=f(x)+g(u)$で、上の分離型の係数空間内の法線空間は
固定した分解$W\oplus U$の下で自然に
\[
 N_{\mathcal S/\mathcal T,F}
 \simeq U\otimes\bigl(\Sym^2W/\langle f_x,f_y,f_z\rangle\bigr)
 \ \oplus\ W\otimes\Sym^2U,
 \qquad \dim N_{\mathcal S/\mathcal T,F}=6+9=15
 \tag{N}\label{eq:split-normal-directions}
\]
と同定される。前者が中心変数に一次の本質的な混合項、
後者が中心変数に二次の混合項を表す。
\end{corollary}
\begin{proof}
全三次式の混合部分は
$\Sym^2W\otimes U\oplus W\otimes\Sym^2U$である。
分離型の接空間の混合部分は、線形shearによる
$\langle f_x,f_y,f_z\rangle\otimes U$と一致する。
滑らかな$f$では三偏導関数が独立なので
その商を取れば上式を得る。
\end{proof}

\begin{proposition}[可換cubicへの忘却と$17=10+7$の接空間]
\label{prop:jacobian-forgetful-tangent}
上の分離点$F=f(x)+g(u)$で、$f$が滑らかな平面三次式、
$g$が重根のない二元三次式とする。$J(F)$を$F$のJacobianイデアルとすると、
\[
  J(F)_3^{\mathrm{quot}}:=\Sym^3(W\oplus U)/J(F)_3
   \simeq J(f)_3^{\mathrm{quot}}
     \oplus\bigl(J(f)_2^{\mathrm{quot}}\otimes J(g)_1^{\mathrm{quot}}\bigr)
     \oplus\bigl(J(f)_1^{\mathrm{quot}}\otimes J(g)_2^{\mathrm{quot}}\bigr)
\]
であり、三項の次元は順に$1,6,3$、合計10である。
ここで$J(f)_d^{\mathrm{quot}}=(\Sym W/(f_x,f_y,f_z))_d$、
$J(g)_d^{\mathrm{quot}}=(\Sym U/(g_u,g_v))_d$とする。
放物型軌道の接空間を割った17次元の変形空間から、
元の可換cubicのJacobian変形空間への自然な線形射
\[
 \Sym^3(W\oplus U)/T_F(\mathcal P\cdot F)
 \longrightarrow J(F)_3^{\mathrm{quot}}
 \tag{J}\label{eq:forgetful-Jacobian}
\]
は全射で、その核は7次元である。
分離型の2次元商に由来する方向の像は1次元、
その余次元15の法線からの像は9次元であり、
法線の核の次元は6である。
\end{proposition}
\begin{proof}
$f$と$g$の勾配は別変数の正則列だから、Jacobian環は
$J(f)^{\mathrm{quot}}\otimes J(g)^{\mathrm{quot}}$に分解し、
Hilbert級数はそれぞれ$(1+t)^3$と$(1+t)^2$となる。
従って次数3には$(3,0),(2,1),(1,2)$の三組だけが寄与し、
次元は$1+3\cdot2+3\cdot1=10$である。
$\mathcal P$の軌道の無限小作用は$F$のスカラー倍と
$X_j\partial_i F$の線形結合である。
Euler等式により$F\in J(F)_3$なので、
\eqref{eq:forgetful-Jacobian}はwell-definedであり全射である。
その次元差は$17-10=7$。
分離型の接空間の純三次方向は$\Sym^3W\oplus\Sym^3U$、
残りは$u f_i,v f_i$なので、Jacobian環への像は
$J(f)_3^{\mathrm{quot}}$の1次元のみである。
さらに系\ref{cor:split-normal-directions}の法線の第一項は
$J(f)_2^{\mathrm{quot}}\otimes J(g)_1^{\mathrm{quot}}$の6方向に写り、
第二項は$J(f)_1^{\mathrm{quot}}\otimes J(g)_2^{\mathrm{quot}}$の3方向に写る。
よって法線からの像は9次元、核は$15-9=6$である。
\end{proof}
\begin{remark}[幾何学的17次元候補との関係]
\label{rem:forgetful-geometry-not-model}
この接空間の$17=10+7$と、幾何学的候補
$(Y,\Lambda,\xi)$の$10+6+1$は整合的である。
実際、元の可換cubicを固定するPoisson方向は
超平面pencilの6方向とPoissonスカラーの1方向を持つ。
しかしこの計算だけでは、量子化後の可換cubic$Y$、平面$\Lambda$、
楕円曲線上の直線束$\xi$を関手的に構成したことにはならない。
\end{remark}

\begin{proposition}[七次元の核の幾何的分解]\label{prop:vertical-seven}
命題\ref{prop:jacobian-forgetful-tangent}の仮定の下で、
$W=\langle x,y,z\rangle$、$U=\langle u,v\rangle$を
\emph{一次形式の空間}とする。分離点$F=f(x)+g(u,v)$で、
\[
 \mathsf T_{F,U}:=\Sym^3(W\oplus U)/T_F(\mathcal P\cdot F),
 \qquad R_F=\Sym(W\oplus U)/(F_x,F_y,F_z,F_u,F_v)
\]
とおく。一次形式の部分空間$U$で指定される平面
$\Lambda=V(U)\simeq\PP^2$の変形空間は
$T_{\Lambda}\operatorname{Gr}(2,W\oplus U)
\simeq\operatorname{Hom}(U,W)$である。
$L\in\operatorname{Hom}(U,W)$の基底表示を
$L(u)=\ell_u(x)$、$L(v)=\ell_v(x)$とすれば、
次の短完全列が成立する：
\[
 0\longrightarrow
 \mathbb C[f]\oplus T_{\Lambda}\operatorname{Gr}(2,W\oplus U)
 \xrightarrow{\ \iota\ } \mathsf T_{F,U}
 \xrightarrow{\ \mathrm{Jac}\ } (R_F)_3
 \longrightarrow0,
 \tag{V}\label{eq:vertical-seven-exact}
\]
\[
 \iota(a,L)=[\,a f+\ell_u(x)g_u(u,v)+\ell_v(x)g_v(u,v)\,].
 \tag{VM}\label{eq:vertical-moving-plane}
\]
左端の一次元方向はPoisson双ベクトルの大きさを変える方向であり、
六次元方向は可換cubicそのものを変えずに指定した平面を動かす方向である。
\end{proposition}
\begin{proof}
$F$のJacobianイデアルの三次部分を変数の双次数に分けると、
\[
 J(F)_3=
 W\langle f_x,f_y,f_z\rangle\ \oplus\
 U\langle f_x,f_y,f_z\rangle\ \oplus\
 W\langle g_u,g_v\rangle\ \oplus\
 U\langle g_u,g_v\rangle.
\]
$f$が滑らかな平面三次式であることから
第一項の次元は$10-\dim(R_f)_3=9$、第二項は6である。
$g$が重根を持たない二元三次式なので
第三項は6、第四項は4であり、全体の次元は25となる。
一方、$\mathcal P$の無限小作用は、
第一項において$\mathrm{SL}(W)$の8次元方向を占め、
第二項ではshearによる全6方向を占め、
第四項では$\mathrm{GL}(U)$による全4方向を占める。
ここで$\mathrm{GL}(U)$が$\Sym^3U$に作用する微分は同型である：
重根のない三点の組の射影的安定化群が有限だからである。
従って放物型軌道の次元は$8+6+4=18$であり、
$W\langle g_u,g_v\rangle$の第三項には寄与しない。
また$\mathbb Cf\subset W\langle f_x,f_y,f_z\rangle$は
$\mathrm{SL}(W)$の8次元接方向と独立である。
実際、滑らかな平面三次式の判別式は$\mathrm{SL}(W)$不変で
$f\mapsto(1+\varepsilon)f$で非零の一次変化を持つ。
従って$J(F)_3/T_F(\mathcal P\cdot F)$は
$\mathbb C[f]\oplus W\langle g_u,g_v\rangle$に一致する。
$g_u,g_v$は二元二次式として一次独立なので、
後者は$\operatorname{Hom}(U,W)$と同一視される。
最後に$\operatorname{Gr}(2,W\oplus U)$の平面変形は
$U\mapsto\langle u+\varepsilon\ell_u(x),
 v+\varepsilon\ell_v(x)\rangle$で記述される。
$g$へこの置換を適用したときの一次変化が
\eqref{eq:vertical-moving-plane}の第二項に一致する。
$F$のJacobian商への射は全射だから、短完全列が得られる。
\end{proof}

\begin{proposition}[PoissonスカラーをPicardパラメータと直接同一視できない]\label{prop:poisson-picard-obstruction}
$X\subset\PP^4$を滑らかな可換cubic、$\Lambda\simeq\PP^2$を
$C=X\cap\Lambda$が滑らかな平面三次曲線となる平面とする。
$X$と$\Lambda$を固定し、非零Poisson構造のスカラー倍
$\lambda\pi$、$\lambda\in\mathbb G_m$だけを変化させる。
この一変数族から$\operatorname{Pic}^0(C)$への任意の
有理写像は定数である。
とくに、同じ$X$と$\Lambda$を保つような
Poissonパラメータ$\lambda$と楕円Picardパラメータ$\xi$の間には、
支配的な有理対応を与える一価の有理写像は存在しない。
\end{proposition}
\begin{proof}
$\operatorname{Pic}^0(C)$は滑らかな射影的種数1曲線である。
$\mathbb G_m\dashrightarrow\operatorname{Pic}^0(C)$という有理写像は、
領域を$\PP^1$に埋め込めば、曲線上の有理写像に関する
正規性とターゲットの固有性により
$\PP^1\to\operatorname{Pic}^0(C)$に延長される。
標数0で$\PP^1$から種数1曲線への非定数射は
Riemann--Hurwitz公式に反する。
従って有理写像は定数である。
この議論は$(X,\Lambda)$を変数とする基底の
一般幾何ファイバー上でも成立する。
\end{proof}
\begin{remark}[復元に必要な非自明な幾何変化]\label{rem:picard-obstruction-scope}
滑らかなcubic threefold上の非零Poisson構造を射影化すると
$\operatorname{Gr}(2,5)$になることは既知である
\cite{LPT,Mayanskiy}。
非零スカラーまで保持するとその錐から零切断を除いた
$\mathbb G_m$束であり、固定cubic上で$6+1$次元を持つ。
一方$(Y,\Lambda,\xi)$の$\xi$は固定$(Y,\Lambda)$上では
楕円曲線$\operatorname{Pic}^0(Y\cap\Lambda)$を動く。
従って$17=10+6+1$という次元一致を、
可換cubicと平面を両方固定したまま
$\mathbb G_m$パラメータを楕円曲線の点へ置き換える
大域的な有理同一視と解釈してはならない。
実際の量子化後の可換復元では、元の$X$と異なる$Y$を選ぶか、
平面やmarkingを非自明に変化させる必要があり得る。
一ブロックHesse模型の$\vartheta(h,t)-t\ne0$は
そのような基底の移動が起こる明示的な例である。
この命題は一般のrotation付き可換復元の存在を否定しない。
また、有限被覆を介する多価の代数的対応の不存在も主張しない。
\end{remark}

\begin{remark}[可換復元への帰結]
\label{rem:split-implication}
分離型のrotation付き可換復元が成立しても、
それだけでは17次元一般点の可換復元は得られない。
線形shearで除去できない本質的な混合項に対する
相対Orlov比較・smooth proper性・rotationの同定が必要である。
ここで求めた余次元は、中心平面を指定する
Poisson中心三次対の\emph{有理商}についての結論であり、
無標識導来圏やKuznetsov成分のモジュライの余次元だと
読み替えてはいけない。
\end{remark}

\section{全混合係数を保つ中心三次式と規格化}
\label{sec:pair-explicit-centers}
中心元を加える操作は同じ周囲の中でも三次対を変えるので、
その中央スカラーを含む規格化を固定する。
以下の有限公式はRainsの中心計算\cite[\S9.1]{EG}と
\href{https://github.com/uedakazushi/nc-cubic/blob/main/papers/support/papers/sections/hesse_mixed_universal_center.tex}
{論文Aの普遍Hesse中心公式}に関連する既知の入力である。
ここでは対の復元に必要な全係数の表示と、関係イデアル内の直接の証明を与える。

\subsection{相対Hesse chartの有限式}
$R=k[u,v]$、$L_i\in R_1$、$M_i\in R_2$、$g\in R_3$として
\[
 F=\frac{x^3+y^3+z^3}{3}+\lambda xyz
   +\frac12(L_1x^2+L_2y^2+L_3z^2)+M_1x+M_2y+M_3z+g
\]
を考える。$\lambda^3\ne-1$なら最高次平面三次式は滑らかである。
任意の滑らかな最高次部分はHesse形へ移せる。
$\lambda=0$でも、$\omega^3=1$, $\omega\ne1$, $c^3=1/3$として
$(x,y,z)=c(X+Y+Z,X+\omega Y+\omega^2Z,X+\omega^2Y+\omega Z)$
とすればHesse係数は2となるので、必要なら$\lambda\ne0$にできる。
一般の二次混合項の交差係数を$C_{12},C_{23},C_{31}\in R_1$とすると、
$x=X-C_{23}/\lambda$, $y=Y-C_{31}/\lambda$, $z=Z-C_{12}/\lambda$
で交差三項を消せる。新たな対角係数は$L_i=\widetilde L_i+2a_i$、
一次係数は$M_i=\widetilde M_i+a_i^2+\widetilde L_i a_i-\lambda a_ja_k$である。
定数は$g$へ吸収する。従ってこのchartは全21混合方向を扱い、
モジュライ上の一意な座標選択を主張しない。

$q=1+\lambda t$と置き、$R$を中心として
\[
 xy-qyx=t(z^2+L_3z+M_3),\quad
 yz-qzy=t(x^2+L_1x+M_1),\quad
 zx-qxz=t(y^2+L_2y+M_2)
\]
で$B_t$を定める。略記を
\begin{align*}
 A&=q^3-t^3,& Q&=q^2+q+1,\\
 K&=q^4+q^3+q^2-(2q+1)t^3,& J&=Q+(q+2)t^3,\\
 H_1&=(q+1)(q^3+q^2+q-t^3),\\
 H_2&=q^3+2q^2+2q+1+(q+1)t^3,\\
 H_3&=q^4+2q^3+2q^2+q-(q+1)t^3
\end{align*}
とする。
\begin{proposition}[可換点で正則なRains規格化]
\label{prop:pair-rains-polynomial}
次の有限式は$D(t)=q(q^3-t^3)(q+1)(1+t^3)\ne0$上で中心正則元となる：
\begin{align*}
 6(G_t-g)={}&(q+1)\{A(x^3+z^3)+(1+t^3)y^3\}\\
 &+q(q+1)(\lambda Q-2t^2)zyx\\
 &+K(L_1x^2+L_3z^2)+JL_2y^2\\
 &-(q+1)t^2(qL_1zy+L_2zx+qL_3yx)\\
 &+(H_1M_1-qt^3L_1^2-qt^2L_2L_3)x\\
 &+(H_2M_2+t^3L_2^2-qt^2L_1L_3)y\\
 &+(H_3M_3-qt^3L_3^2-qt^2L_1L_2)z.
\end{align*}
$D(0)=2$、$G_0=F$であり、$B_t/(G_t)$はHilbert級数
$(1-s^3)/(1-s)^5$を持つ平坦な次数付き族である。
\end{proposition}
\begin{proof}
$L_i,M_i$を独立中心変数、重みをそれぞれ1,2として扱う。
重み4の自由語339個に対する二次関係の全倍数162本を
$\mathbb Q(q,t)$上で並べるとrank144となり、表示式の三交換子はその行空間に入る。
\texttt{verify\_central\_pair\_centers.py}がこの関係行列と式を再生成して厳密還元する。
係数を数値へ代入した補間ではない。
$q+1\ne0$では関係式を$(q+1)/2$で割ると、
量子係数$h=2t/(q+1)$、最高次Hesse係数$\lambda$の対称化potentialになる。
補題\ref{lem:pbw}の一様PBWとCalabi--Yau証明をこの係数で適用できる。
この平坦性により分数体上の中心性は$D\ne0$全体へ降下する。
次数3の関係行列の17次小行列式は$-q^5(q^3-t^3)^2$であり、
整列三次語の基底で最高次成分は
\[
 (q+1)(1+t^3)(x^3+z^3)+(q+1)(q^3-t^3)y^3
 +(q+1)(\lambda Q-2t^2)xyz
\]
となる。これは非零なので、整域性から$G_t$は正則である。
商のassociated gradedは最高次中心元による商と$R$のtensor積となり、
その基底の持上げで自由性とHilbert級数を得る。
\end{proof}
\begin{remark}
EG掲載のRains式の$y$係数の$q^3b_2$は$q^3$と読む必要がある。
元のMAGMA資料\cite{RainsMagma}と上の独立な関係行列で照合する。
中心性自体は既知入力の具体化であり、全特殊パラメータのtails正則性は含まない。
\end{remark}

\subsection{Fermatの六二次係数に対する自由代数恒等式}
$(x_1,x_2,x_3)=(x,y,z)$とし、添字$(i,j,k)$を巡回させる。
$L_i,e_i$の重みは1、$m_i$の重みは2とする。
\[
 F=\frac13\sum_i x_i^3+\frac12\sum_iL_ix_i^2
       +\sum_i e_i\Symm(x_jx_k)+\sum_i m_ix_i,
 \qquad g_i=x_i^2+L_ix_i+e_jx_k+e_kx_j+m_i.
\]
$[x,y]=hg_3$, $[y,z]=hg_1$, $[z,x]=hg_2$で関係を定める。
\begin{lemma}[全Artin商での形式PBW]
全ての$N\ge1$で、$h^N=0$上のこの代数は順序付き語$x^ay^bz^c$を
係数環上の基底に持つ。
\end{lemma}
\begin{proof}
$yx\mapsto xy-hg_3$, $zx\mapsto xz+hg_2$, $zy\mapsto yz-hg_1$を用いる。
残る$h$次数と転倒数の辞書式順序で停止する。
唯一のoverlap $zyx$の二経路の差は
$h([x,g_1]+[y,g_2]+[z,g_3])=0$である。
最後は自由代数恒等式で、交差係数の交換子が対で消える。
停止・局所合流性から基底を得て、各$N$の商写像とも両立する。
有限の非零$h$における整列基底を仮定していない。
\end{proof}
\begin{theorem}[四次補正で止まる全係数中心式]
\label{thm:pair-full-fermat-center}
$\Omega_h=\sum x_i^3/3-h^2\sum_{\sigma\in S_3}\sigma(xyz)/9$とする。
$L_i=0$のとき
\[
 Z_h=\Omega_h+\sum_i e_i\Symm(x_jx_k)+\sum_i m_ix_i
       +\frac{h^2}{6}\sum_i e_i^2x_i
       +\frac{h^4}{18}\sum_i(e_ix_i^2-e_je_kx_i)
\]
は全有限$h$・全係数で中心的である。一般の$L_i$に対しては
\begin{align*}
 X_i&=x_i+L_i/2,\qquad
 n_i=m_i-L_i^2/4-(e_jL_k+e_kL_j)/2,\\
 c_{\rm cl}&=\sum_iL_i^3/12+\sum_ie_iL_jL_k/4-\sum_im_iL_i/2,\\
 W_h(x;L,e,m)&=Z_h(X;e,n)+c_{\rm cl}
\end{align*}
が中心元で、$W_0=F$となる。
\end{theorem}
\begin{proof}
$L=0$の自由代数で$r=[x,y]-hg_3$, $s=[y,z]-hg_1$, $t=[z,x]-hg_2$とする。
次の四つの式を定める：
\begin{align*}
 C_0={}&-(ry^2+yry+y^2r)/3+(tz^2+ztz+z^2t)/3+m_3t-m_2r\\
       &-e_3(rx+xr)/2+e_2(tx+xt)/2+e_1(yt-rz+[x,s]/2),\\
 C_1={}&([yz+zy,s]+e_3[t,z]+e_2[r,y])/6,\\
 C_2={}&(r[x,z]-sx^2-t(xy+yx)-xrz+xsx+xty\\
       &\hspace{6em}+2ytx-yxt+zxr)/18,\\
 C_3={}&([r,y^2]+[t,z^2]+e_1[s,x])/18.
\end{align*}
自由代数内で$[Z_h,x]=C_0+hC_1+h^2C_2+h^3C_3$である。
右辺は41個の関係式倍であり、有理数係数の展開で直接検証できる。
巡回置換が残る二つの交換子を与える。
平行移動については$X_i^2+e_jX_k+e_kX_j+n_i=g_i$であり、
$\sum X_i^3/3+\sum e_i\Symm(X_jX_k)+\sum n_iX_i+c_{\rm cl}=F$は
自由代数の恒等式なので、全ての主張が従う。
\end{proof}
\begin{remark}[規格化と不可欠な四次項]
$L=0$では$\det\operatorname{Hess}F=8xyz-2\sum e_i^2x_i+2e_1e_2e_3$である。
従って$Z_h=\Symm(F-h^2\det\operatorname{Hess}F/12)+h^2e_1e_2e_3/6
+h^4\sum(e_ix_i^2-e_je_kx_i)/18$。
最後の項は省けない。$F=\sum x_i^3/3+rxy$では省いた候補の$x$との交換子は
$-h^5(ry^2z+r^2xz)/9\pmod{h^6}$で、形式PBWにより非零である。
中央平行移動$x\mapsto x+a$に対し
$L_i'=L_i+2a_i$, $m_i'=m_i+L_ia_i+a_i^2+e_ja_k+e_ka_j$とすると
\[
 W_h(x+a;L,e,m)=W_h(x;L',e,m')+F(a).
\]
対角Rains式（$\lambda=0,q=1,t=h$）とpreferred式の差は
\[
 G_h^{\rm R}-W_h(x;L,0,m)
 =\frac{h^2}{12}L_1L_2L_3+
   \frac{h^3}{6}(L_1m_1-L_2m_2+L_3m_3).
\]
この中央三次差を捨てると、同じ周囲の中心対を誤同定する。
\end{remark}
\begin{proposition}[有限微分性と形式正則商]
順序付きPBW固定化の積$\mu_h$の$h^r$係数は、係数$L,e,m$を微分しない
$x,y,z$の双微分作用素で、入力の全微分次数は高々$2r$である。
また中心商は各$h^N=0$上で係数環上自由となる。
\end{proposition}
\begin{proof}
Barmeier--Wang\cite[Theorem 10.3, Lemma 10.9]{PairBW}の還元グラフ表示では
$h^r$に寄与する頂点数は$r$であり、各頂点から入力への微分は高々二つ。
非零補正はactive三変数の三対だけなので、中央九係数の微分は現れない。
PBWの合流性によりこの積は指定順序の実際の積である。
中心性と結合性は同じ固定化の曲率付きMaurer--Cartan方程式を与える。
商の自由性は、$F$が$x$のmonic三次の定数倍であることから直接示せる。
$x$次数二以下の単項式空間を$V$とすると、$(a,v)\mapsto Fa+v$は線形同型。
その変形$(a,v)\mapsto W_h*a+v$は元の同型と$1+hT$の合成であり、
$h^N=0$では有限幾何級数で逆を持つ。
この分裂は任意の中央係数代入後にも保たれ、$g(u,v)$を加えても同じである。
\end{proof}

\begin{corollary}[可換点を含む一様な有限係数量子化]
滑らかな最高次平面三次式$f_3$を固定すると、$\Delta(0)\ne0$なる一変数多項式
$\Delta(h)$が存在し、その非零範囲では全21混合係数を同時に保つ中心元$G_h$を代数的に取れる。
$G_0=F$で、中心商は次数ごとに平坦である。
さらに$h\ne0$なら$\dim Z(B_h)_1=2$, $\dim Z(B_h)_3=5$とできる。
\end{corollary}
\begin{proof}
Hesse座標の選択後に$t=2h/(2-\lambda h)$, $q=(2+\lambda h)/(2-\lambda h)$を用いると、
命題\ref{prop:pair-rains-polynomial}が対称化Jacobian関係の中心元を与える。
その分母と$D(t)$の分子は$h=0$で非零で、混合係数に依存しない。
中心の低次数は交換子行列を$h$で割って調べる。
$h=0$でこれは$f_3$のJacobian Poisson中心を計算し、次数1,2,3のrankは3,6,9である。
その非零minorをさらに反転する。三次の一般rankは中心元の存在から9を超えず、
従って三次核はrank1である。
低次補正の未知数9個の行列は、active次数の順に、上の次数2・1の
divided-commutator行列を対角に持つ三角blockを含む。
同じminorを反転すれば、全混合係数に関する唯一の補正を代数的に解ける。
これら有限個の因子の積を$\Delta(h)$とする。
相対PBWで最高次中心元を引く帰納法から中心の次数1,3の等式を得る。
$h=0$ではdivided-commutatorの核はPoisson中心であり、実際の可換中心は
次数1で5、次数3で35に跳ぶ。この二つを同一視していない。
\end{proof}

\section{37次元の周囲成分とその一般中心}
\label{sec:pair-generic-component}
一般のPoisson成分分類の正本は
\href{https://github.com/uedakazushi/nc-cubic/blob/main/papers/paper_a.tex}{論文A}である。
本節で用いる入力は、五変数の標準divergence-free liftの錐
$C_{113}^{\rm aff}$が37次元の被約既約成分であることと、
Pym\cite[Proposition 2.9, Theorem 2.10]{Pym}の一次項が恒等な形式量子化比較である。
ここではこの入力から中心三次対の一般点を具体的に記述する。

$V=k^5$とし、$V^{\otimes2}=\wedge^2V\oplus\Sym^2V$で規格化して
\[
 R_\Pi=\langle a\otimes b-b\otimes a-\Symm(\Pi(a,b)):a,b\in V\rangle,
 \qquad B_\Pi=T(V)/(R_\Pi)
\]
と置く。
\begin{proposition}[graph成分の代数的表示]
$\Pi\mapsto R_\Pi$は$C_{113}^{\rm aff}$を二次関係Grassmannianの
標準アフィンchartへ閉埋込みする。その像$Q_{113}$は37次元で、
可換点の規格化CY変形の被約芽において対応する既約成分をなす。
\end{proposition}
\begin{proof}
$R_\Pi$は$\wedge^2V\to\Sym^2V$のgraphなので$\Pi$を一意に回復でき、
標準liftが線形であることから閉埋込みとなる。
最高次平面三次制限が滑らかな相対Jacobian点は稠密であり、
その点では補題\ref{lem:pbw}と\ref{lem:CY}がKoszul CY代数を与える。
CY変形の条件はこの変形芽では閉条件なので、閉包の芽も含まれる。
それを含む被約既約成分の芽を$D$とする。
一次項が恒等な形式比較の下で$D$の接錐の被約台はあるPoisson既約成分$C_D$である。
graph像は斉次錐なので
$C_{113}^{\rm aff}=TC_0Q_{113}\subset TC_0D=C_D$となる。
Poisson成分の極大性から両者は等しく、次元も37である。
従って$D$と$Q_{113}$の被約な既約芽は等しい。
これはKontsevich級数自体がgraph式と一致するという主張ではない。
\end{proof}

\begin{lemma}[相対Jacobian括弧の全Poisson中心]
$R=k[u,v]$とし、$F\in R[x,y,z]$の最高active次数部分$f_3$が既約とする。
$R$を中心、$\{x,y\}=F_z$などと定めると$Z_{\rm Pois}(R[x,y,z])=R[F]$である。
\end{lemma}
\begin{proof}
中心元の最高次部分$g_m$は$df_3\wedge dg_m=0$を満たす。
Euler場との縮約から$3f_3dg_m=mg_mdf_3$。
$f_3$は素元で、少なくとも一つの非零偏微分を割らないので$f_3\mid g_m$である。
これを繰り返すと$m=3j$、$g_m=cf_3^j$ ($c\in R$)となる。
$g-cF^j$へ次数帰納法を適用すればよい。逆包含はCasimir性による。
\end{proof}
\begin{theorem}[量子成分の関数体上の中心]
\label{thm:pair-generic-center}
$K=k(Q_{113})$とし、一般ファイバーを$B_\eta$とすると
\[
 Z(B_\eta)=K[a,b,c],\qquad\deg a=\deg b=1,\quad\deg c=3.
\]
$B_\eta$は整域で、中心の次数1,2,3の次元は2,3,5となる。
全次数の中心が一つのZariski開部分で一定であるとは主張しない。
\end{theorem}
\begin{proof}
中心変数を厳密に保つ相対Jacobian量子化と中心元$G_\hbar\equiv F\pmod\hbar$を取る。
次数$d$の形式中心元を$\hbar$で還元すると前補題のPoisson中心に入り、
$u,v,G_\hbar$の同じ重みの多項式を引けば差は$\hbar$で割れる。
反復と完備性から
\[
 Z(B_\hbar)_d=\bigoplus_{i+j+3l=d}k[[\hbar]]u^iv^jG_\hbar^l.
\]
独立性は関係式の最小$\hbar$次数を取ることで従う。

この計算を成分の一般点へ戻す。
標準次数付き整域$A$の頂点完備化から、支配的パラメータ表示$T\to\Spec A$を用いて
\[
 \widehat A\longrightarrow k(T)[[\hbar]],\quad
 \sum a_n\longmapsto\sum\hbar^n\phi^*a_n
\]
を作ると、最初の非零斉次係数が消えないので単射である。
これをPoisson錐と形式量子化比較に適用し、$K$を含む体
$L=k(T)((\hbar))$を得る。$T$はpencilの6座標と三次式の31係数で取れる。
各次数の中心は有限自由空間の交換子行列の核なので体拡大と可換し、
$Z(B_\eta)_d\otimes_K L=L[u,v,G_\hbar]_d$となる。
$Z(B_\eta)_1$の基底$a,b$を取り、$c\in Z(B_\eta)_3\setminus\Sym^3\langle a,b\rangle$を選ぶと、
$L$上では$c=\lambda G_\hbar+r(u,v)$ ($\lambda\ne0$)である。
従って$K[a,b,c]\to Z(B_\eta)$は忠実平坦な体拡大後に同型、ゆえに元から同型である。
整域性も形式PBW量子化への埋込みから従う。
\end{proof}
\begin{corollary}[支配的incidenceと有限次数の良い開部分]
$Q_{113}$の次数三中心直線のincidenceには一般ファイバー$\PP^4_K$を持つ
41次元の支配的既約成分$I_{113}$が一意に存在する。
中心の次数1,3とHilbert関数の次数4以下が一般値を取る開部分では、
これはrank5の中心核束の$\PP^4$束である。
一般の有理商の次元は周囲で13、対で17、忘却の一般ファイバーは4次元である。
\end{corollary}
\begin{proof}
中心の有限行列の非零minorを反転するとrank5の核束を得る。
その射影束は既約で一般ファイバー$\PP^4_K$の閉包を与え、支配的成分は他にない。
相対Jacobianパラメータ族の6+35=41次元の像がこれを支配する。
六分岐値を持つ対角族の周囲は射影stabilizerが有限であることを
第\ref{sec:pair-diagonal-classification}節で証明する。
従って$\PGL_5$の一般軌道は24次元である。
Rosenlichtの定理により37−24=13、41−24=17を得て、忘却の一般ファイバーは4次元となる。
この計算は商多様体の有理性や一般stabilizerの自明性を含まない。
\end{proof}

\section{対角混合族の完全な同型条件と反対代数}
\label{sec:pair-diagonal-classification}
$D=x_1^3+x_2^3+x_3^3$、$T$を六つの混合語の和とし、
$s=(u,v)^t$, $q_i\in k[u,v]_2$, $f\in k[u,v]_3$に対して
\begin{gather*}
 B(h,q)=\frac{k[u,v]\langle x_1,x_2,x_3\rangle}
 {([x_j,x_k]-h(x_i^2+q_i))_{(i,j,k)\text{ cyclic}}},\\
 W_{h,q,f}=D/3-h^2T/9+\sum_iq_ix_i+f,
 \quad A(h,q,f)=B(h,q)/(W_{h,q,f})
\end{gather*}
と置く。以下$h\ne0$とする。境界の滑らかさは同型判定には不要である。

\begin{lemma}
$Z(A)_1=ku\oplus kv$であり、$\Sym^2(A_1)\to A_2$は同型である。
\end{lemma}
\begin{proof}
二次関係は$\wedge^2A_1$から$\Sym^2A_1$へのgraphで、三次商は次数二を変えない。
従って対称二次式は$A_2$の基底である。
一次元$\sum a_ix_i+bu+cv$の交換子を取ると、独立な$x_i^2$の係数と$h\ne0$から
$a_i=0$を得る。
\end{proof}
\begin{theorem}[対角族の同型の必要十分条件]
\label{thm:pair-diagonal-isomorphism}
$h,h'\ne0$について$A(h',q',f')\simeq_{\gr} A(h,q,f)$であることは、
$p\in S_3$, $N\in\GL_2$, $\lambda_i\ne0$, $\kappa\ne0$が存在して
\begin{gather*}
 \lambda_1^3=\lambda_2^3=\lambda_3^3=\kappa,\qquad
 h'=\operatorname{sgn}(p)h\lambda_1\lambda_2\lambda_3/\kappa,\\
 q_i'(Ns)=\lambda_i^2q_{p(i)}(s),\qquad f'(Ns)=\kappa f(s)
\end{gather*}
を満たすことと同値である。全ての同型は
$s'\mapsto Ns$, $x_i'\mapsto\lambda_i x_{p(i)}$であり、特に$h^6$は不変量である。
\end{theorem}
\begin{proof}
二次包絡の復元により、同型は周囲へ持ち上がる。
中心一次部分から$s'\mapsto Ns$, $x'\mapsto Mx+\ell(s)$であり、
関係式を対称二次式として比較すると
\[
 h\sum_a\operatorname{cof}(M)_{ia}(x_a^2+q_a(s))
 =h'\bigl((Mx)_i+\ell_i(s)\bigr)^2+h'q_i'(Ns).
\]
$x_as_b$の係数から$\ell_i=0$、$x_ax_b$ ($a\ne b$)の係数から
$M_{ia}M_{ib}=0$となる。可逆性により$M$はmonomialである。
$x_{p(i)}^2$の係数を比較して
$h'\lambda_i^2=\operatorname{sgn}(p)h\lambda_j\lambda_k$を得て、
$\lambda_i$を掛ければ三つの立方は等しい。
中央二次部分の比較が$q_i$の条件を与える。
$L=\prod\lambda_i$と置けば$L^3=\kappa^3$, $(h')^2L=\kappa h^2$なので、
$W'$の像は$\kappa(D/3-h^2T/9+\sum q_ix_i)+f'(Ns)$となる。
最高active三次部分の比較から比例係数は$\kappa$であり、$f$の条件を得る。
この最高次元は$h^6=-27$でも非零である。
実際、三変数代数の可換化は$k[x,y,z]/(x^2,y^2,z^2)$で、
$\Omega_h$の像は$-2h^2xyz/3\ne0$である。
逆に表示条件を満たす可逆線形写像は全ての関係を保つので同型を与える。
$h'$の式の六乗が最後の不変量を与える。
\end{proof}

\begin{proposition}[軌道次元と自己同型]
$q_1q_2q_3$の六零点が相異なる範囲で、周囲の対角族の同型類は6次元、
中心三次対の同型類は10次元である。全37次元周囲族と41次元対の中で、
それぞれの$\PGL_5$飽和像の余次元は7となる。
さらに六点集合の$\PGL_2$ stabilizerが自明で$f\ne0$なら
$\Aut_{\gr}A=k^\times$、$A\not\simeq_{\gr} A^{\rm op}$である。
\end{proposition}
\begin{proof}
連続な変換パラメータは$(c,N)\in\mathbb G_m\times\GL_2$の5次元であり、
全生成元の同時スカラーという1次元核を持つ。
六点のstabilizerは有限で、各候補の射影$N$を固定すると$q_i$の条件は
$N$のスカラーと$c$の比を有限個に制限する。
従って軌道は4次元で、周囲の10係数・対の14係数からそれぞれ4を引く。
一般の$\PGL_5$軌道は24次元なので飽和像は30次元・34次元であり、余次元7を得る。
stabilizerが自明なら$N=tI$, $p=1$で、$\lambda_i^2=t^2$と立方の条件から
全ての$\lambda_i$は同じ$\lambda=\pm t$である。
$f\ne0$が$\lambda^3=t^3$を強制し、自己同型は同時スカラーのみとなる。
反対代数は$A(-h,q,f)$であるが、同じ条件は$h'=h$を強制して矛盾する。
\end{proof}
\begin{corollary}[有理係数の非自己反対例と標識付き導来圏]
\label{cor:pair-not-opposite}
$q_1=uv$, $q_2=(u-v)(u-2v)$, $q_3=(u-4v)(u-9v)$とする。
六点集合のstabilizerは自明である。
$h=\pm1$について、両方のtailsが正則となる$f$の非空開集合が存在し、
その中の$f\ne0$では$A$と$A^{\rm op}$の間に構造対象をline objectのshiftへ送る
導来同値は存在しない。有理係数のそのような$f$はZariski稠密に存在する。
\end{corollary}
\begin{proof}
$\{0,\infty,1,2,4,9\}$の最初の三点の像を選ぶ120通りの有理射影行列を列挙し、
全六点を保つものを照合すると恒等変換だけである。
\texttt{verify\_central\_pair\_orbits.py}はこの厳密列挙を行う。
$h=\pm1$の境界の位数は補題\ref{CSBGlem:order-six}により6であり、
定理\ref{CSBGthm:global-open}を両周囲に適用する。
全ての$f$で相対PBW・中心正則商の自由性があり、最高次中心元は非零正則なので全仮定が揃う。
二つの非空開集合は交わり、$f=0$を除いても非空である。
Serre関手・偶数Veroneseによる復元定理と前命題から導来同値の不存在を得る。
有理点の稠密性は$\mathbb Q^4\subset\mathbb C^4$のZariski稠密性による。
構造対象を忘れた全ての導来同値についての主張ではない。
\end{proof}

\section{本質的混合項の量子中心と相対Picardの制約}
\label{sec:mixed-quantum-center}
分離型を越える量子化では中心変数一次の補正が最初の課題となる。
以下の中心元表示はEtingof--Ginzburg \cite[\S9]{EG} のRains型計算および
本研究の\href{https://github.com/uedakazushi/nc-cubic2/blob/main/papers/pdf-only/ambient/central_pair_formal_moduli.pdf}{既存の中心三次対技術稿}
第14節にある結果のFermat規格化であり、新しい中心公式とは主張しない。
従来の独立計算は有限体での次数四イデアル所属検査に留まったが、
現在は論文Aの
\texttt{papers/support/papers/sections/fermat\_mixed\_universal\_center.tex}
で全係数の普遍的な中心性を厳密に証明した。

\subsection{全15混合方向の中心三次元}
$f=(x^3+y^3+z^3)/3$、$U=\langle u,v\rangle$とし、
$A,B,C\in U$、$D,E,H\in\Sym^2U$、$g\in\Sym^3U$を任意に取る。
\begin{equation}\label{eq:full-mixed-fermat}
 F=f+Axy+Byz+Czx+Dx+Ey+Hz+g
\end{equation}
とおき、$u,v$を中心変数として
\begin{equation}\label{eq:full-mixed-brackets}
 [x,y]=hF_z,\quad[y,z]=hF_x,\quad[z,x]=hF_y
\end{equation}
の右辺を完全対称化して読む。$T=xyz+xzy+yxz+yzx+zxy+zyx$、
$\operatorname{sym}(xy)=(xy+yx)/2$とする。

\begin{remark}[既存原稿の全混合中心元と検算状態]\label{rem:full-mixed-center}
既存の技術稿の中心元の表示は、現在の規約で
\begin{align}\label{eq:full-mixed-center}
\Omega_h={}&\frac{x^3+y^3+z^3}{3}-\frac{h^2}{9}T\notag\\
&+A\left(\operatorname{sym}(xy)+\frac{h^4}{18}z^2\right)
 +B\left(\operatorname{sym}(yz)+\frac{h^4}{18}x^2\right)
 +C\left(\operatorname{sym}(zx)+\frac{h^4}{18}y^2\right)\notag\\
&+\left(D+\frac{h^2B^2}{6}-\frac{h^4AC}{18}\right)x
 +\left(E+\frac{h^2C^2}{6}-\frac{h^4AB}{18}\right)y\notag\\
&+\left(H+\frac{h^2A^2}{6}-\frac{h^4BC}{18}\right)z+g(u,v)
\end{align}
となり、$h=0$で$F$に特殊化する。
この有限多項式は現在、論文Aで普遍的な中心元として証明済みである。
\texttt{computations/verify\_universal\_fermat\_mixed\_center.py}
では$A,B,C,D,E,H$を独立な中心不定元とし、
$\hbar$も変数のまま$\mathbb Q(\hbar)$上の次数四関係式行列を厳密に還元する。
行列の列数339、関係式の倍数162、一般階数と$\hbar=0$での階数はいずれも144、
全三交換子の剰余は零である。これにより原点を含む開集合での厳密な中心性が従う。
新たな \texttt{computations/verify\_mixed\_center.py} では
$p=1009$、$h=0,1,2,5,13$、各$h$に3組ずつ混合係数を代入して、
次数四の自由語625個における二次関係式の生成部分空間
（階数555）に全ての交換子が属することを確認した。
以下の有限特殊化は追加の検算であり、現在の普遍的証明の根拠ではない。
$A\ne0$だけの模型で$h^2A^2z/6$を省くと
$h=1$の次数四検査で中心性が失われることも確認した。
\end{remark}

\subsection{一次中心持上げの一意性}
\begin{proposition}\label{prop:mixed-center-jet-unique}
$S_h=\CC\langle x,y,z\rangle/
([x,y]-hz^2,[y,z]-hx^2,[z,x]-hy^2)$とし、
$Z(S_h)_2=0$を仮定する。
$R=\CC[u,v]/(u,v)^2$上の\eqref{eq:full-mixed-brackets}の
量子化族$\mathcal B_R$が次数ごとに$R$-自由なら、
固定された$S_h$の次数三の中心元$\Omega_{h,0}$への
次数三の中心持上げ$\Omega_R\in Z(\mathcal B_R)_3$は、
存在すれば一意である。
\end{proposition}
\begin{proof}
二つの持上げの差$\Delta$は$(u,v)(\mathcal B_R)_3$に属し、
次数別自由性と平方零条件により
$\Delta=uQ+vR'$、$Q,R'\in(S_h)_2$と一意に書ける。
$\Delta$の中心性を各生成元$a=x,y,z$との交換子に適用すると
$u[Q,a]+v[R',a]=0$を得る。
$u,v$が独立な平方零係数なので、
$Q,R'\in Z(S_h)_2=0$。従って$\Delta=0$。
\end{proof}
\begin{remark}[一次混合6方向の明示的持上げ]
\eqref{eq:full-mixed-center}を$(u,v)^2=0$に制限すると
\[
\begin{aligned}
 \Omega_h^{(1)}
 ={}&\frac{x^3+y^3+z^3}{3}-\frac{h^2}{9}T
 +A\left(\operatorname{sym}(xy)+\frac{h^4}{18}z^2\right)\\
 &+B\left(\operatorname{sym}(yz)+\frac{h^4}{18}x^2\right)
 +C\left(\operatorname{sym}(zx)+\frac{h^4}{18}y^2\right).
\end{aligned}
\]
これが全ての一次混合6方向の持上げの表示である。
命題\ref{prop:mixed-center-jet-unique}の一意性自体は、
既存の中心公式を仮定しない独立した証明である。
$h=0$では$Z(S_0)_2\ne0$なので適用しない。
\end{remark}

\subsection{相対Picardに対するスカラー方向の障害}
\begin{proposition}\label{prop:relative-scalar-picard}
$X\subset\PP^4$を滑らかなcubic threefoldとし、
$T_0\subset\operatorname{Gr}(2,H^0(\PP^4,\OO(1)))$を、
指定する平面$\Lambda_U=V(U)$との切断$C_U=X\cap\Lambda_U$が
滑らかな楕円曲線になる非空開部分とする。
$\mathscr C\to T_0$を楕円曲線族、
$\mathscr P=\operatorname{Pic}^0(\mathscr C/T_0)$を相対Picardとする。
非零Poisson錐の$\mathbb G_m$-束
$\mathscr L^\times\to T_0$から$\mathscr P$への
$T_0$-上の任意の有理写像は、稠密開部分において
$\mathscr L^\times\to T_0$を通じて有理的に因子化する。
すなわちPoissonスカラー方向には定数である。
幾何学的一般ファイバーで連結な有限\emph{étale}被覆を許しても同様である。
\end{proposition}
\begin{proof}
$K=\CC(T_0)$とする。$\mathscr L^\times$の一般ファイバーは
$\mathbb G_{m,K}$である。
$\overline K$へ移すと、与えられた有理写像は
$\mathbb G_m\dashrightarrow E$、$E$は種数1曲線となる。
固有性によって$\PP^1\to E$に延長されるが、
Riemann--Hurwitz公式から定数である。
従ってもとの写像は一般ファイバーに沿って定数であり、
$T_0$上の有理sectionを通じて因子化する。
また代数閉体上の$\mathbb G_m$の連結有限étale被覆は
$t\mapsto t^n$型で、その源も$\mathbb G_m$である。
同じ議論が適用できる。
\end{proof}
\begin{remark}[復元への含意と限界]
元の$X$とPoisson pencil $U$を幾何側でも固定したまま、
非零スカラーを楕円Picardの非自明な点へ送る大域有理模型は存在しない。
しかし量子化後に別の可換cubic $Y$や平面が動く比較、
分岐被覆を介した多価対応、形式局所同値の不存在は結論しない。
一般の非零パラメータでのrotation付き可換復元に必要な比較は未完成である。
有限Artin底での線束生成残余pairと相対MFの比較は、
論文CのArtin残余比較の節で扱う。
\end{remark}


\section{正則正規超曲面とAS正則$\mathbb Z$-代数の周期性}\label{sec:moduli-ambient-periodicity}
本節では、通常の次数付き表示をあらかじめ仮定しないAS正則$\mathbb Z$-代数を扱う。
これは本稿の中心三次対のモジュライを一般のhelix型ambientへ拡張する際の構造的制約であり、
本稿で直接構成する次数付きambientの1周期性とは区別する。
以下の周期性はAS双対性と正規族の次数から従うものであり，
全ての抽象的な非可換cubic threefoldが超曲面表示を持つという主張ではない。

局所有限な正の連結 $\mathbb Z$-代数
$\mathcal B=\bigoplus_{i,j\in\mathbb Z}\mathcal B_{ij}$ に対し，
$\mathcal B_{ii}=\CC$，$\mathcal B_{ij}=0$（$j<i$）とする。
$\mathcal B(r)_{ij}=\mathcal B_{i+r,j+r}$ と置き，
添字を保つ同型 $\mathcal B\simeq\mathcal B(r)$ があるとき
$r$ \emph{周期的}という。
通常の connected graded algebra $D=\bigoplus_{n\geq0}D_n$ に
付随する $\check D_{ij}=D_{j-i}$ は1周期的である。
逆に1周期同型を選ぶと，それを用いて $\mathcal B_{0n}$ の積を
定めることで $\mathcal B\simeq\check D$ と復元できる。
この同値では周期同型そのものを選ぶことに注意する。

次数 $d>0$ の\emph{正則正規族}とは，
$g_i\in\mathcal B_{i,i+d}$ の族で，すべての $i\leq j$ について
\begin{equation}\label{eq:moduli-ambient-normal-family}
 g_i\mathcal B_{i+d,j+d}
 =\mathcal B_{ij}g_j\subset\mathcal B_{i,j+d}
\end{equation}
を満たし，各 $g_i$ の左右乗法が対応する成分上で単射となるものをいう。
正規族の定義については
Chan--Nyman~\cite[Definition~8.6]{PChanNymanNormalFamily}も参照されたい。
ここでは中心性を仮定しない。

\begin{theorem}[正則正規超曲面の ambient の周期]\label{thm:moduli-ambient-gcd}
$\mathcal B$ は二次Koszul AS正則 $\mathbb Z$-代数で，
左右のAS条件を満たし，各単純加群は有限生成射影加群による
有限の最小線型解消を持つとする。
そのAS次元とGorenstein parameterが同じ $N$ であり，
次数 $d$ の正則正規族が存在するならば，最小正周期 $p$ は
\begin{equation}
 p\mid\gcd(N,d)
\end{equation}
を満たす。特に $\gcd(N,d)=1$ なら
$\mathcal B\simeq\mathcal B(1)$ である。
\end{theorem}

\begin{proof}
まず正規族から $d$ 周期性を構成する。
各 $a\in\mathcal B_{ij}$ に対して
\eqref{eq:moduli-ambient-normal-family}と左右正則性から，
\begin{equation}
 ag_j=g_i\sigma_{ij}(a),\qquad
 \sigma_{ij}(a)\in\mathcal B_{i+d,j+d}
\end{equation}
を満たす元が一意に存在し，$\sigma_{ij}$ は線型同型となる。
$a\in\mathcal B_{ij}$，$b\in\mathcal B_{j\ell}$ について
\begin{equation}
 (ab)g_\ell
 =a(g_j\sigma_{j\ell}(b))
 =g_i\sigma_{ij}(a)\sigma_{j\ell}(b)
\end{equation}
だから，$g_i$ を消去して
$\sigma_{i\ell}(ab)=\sigma_{ij}(a)\sigma_{j\ell}(b)$ を得る。
従って
\begin{equation}\label{eq:moduli-ambient-sigma}
 \sigma:\mathcal B\xrightarrow{\sim}\mathcal B(d)
\end{equation}
は $\mathbb Z$-代数同型である。
さらに$a=g_i$として同じ消去を用いれば
$\sigma_{i,i+d}(g_i)=g_{i+d}$であり、$d$周期同型は正規族を保つ。

次にKoszul双対 $E=\mathcal B^!$ を考える。
有限線型解消と左右AS双対性から，$E$ はindex $N$ の
局所Frobenius $\mathbb Z$-代数となる
（Bondal--Polishchuk~\cite[Proposition~4.1, \S5]{PBondalPolishchukHelices}）。
具体的には，最高次数の同定
$\lambda_i:E_{i,i+N}\xrightarrow{\sim}\CC$ を選ぶと，
\begin{equation}
 E_{ij}\otimes E_{j,i+N}\longrightarrow E_{i,i+N}
 \xrightarrow{\lambda_i}\CC
\end{equation}
が各相補次数で完全対合である。
この対合から，任意の $a\in E_{ij}$ に対して
\begin{equation}
 \lambda_i(ab)=\lambda_j(b\nu_{ij}(a))
 \quad(b\in E_{j,i+N})
\end{equation}
となる線型同型 $\nu_{ij}:E_{ij}\to E_{i+N,j+N}$ が定まる。
$a\in E_{ij},c\in E_{j\ell},b\in E_{\ell,i+N}$ に対して
\begin{equation}
 \lambda_i((ac)b)
 =\lambda_j((cb)\nu(a))
 =\lambda_\ell(b\nu(a)\nu(c))
\end{equation}
であるから，対合の完全性により
$\nu(ac)=\nu(a)\nu(c)$ となる。
従って $\nu:E\xrightarrow{\sim}E(N)$ は $\mathbb Z$-代数同型である。
次数一部分で逆双対を取り，二次関係を保つことを用いると，
\begin{equation}\label{eq:moduli-ambient-rho}
 \rho:\mathcal B\xrightarrow{\sim}\mathcal B(N)
\end{equation}
が得られる。

周期の集合
$\operatorname{Per}(\mathcal B)=
\{r\in\mathbb Z:\mathcal B\simeq\mathcal B(r)\}$
は $\mathbb Z$ の部分群である。
\eqref{eq:moduli-ambient-sigma}と\eqref{eq:moduli-ambient-rho}から
$d,N\in\operatorname{Per}(\mathcal B)$ なので，
$\gcd(N,d)$ も周期となる。
二つの同型が可換であることは必要ない。
\end{proof}

\begin{corollary}[三次三様体の ambient]\label{cor:moduli-cubic-ambient-periodic}
$\mathcal B$ を標準 $(1,5,10,10,5,1)$ 型の
5次元Koszul AS正則 $\mathbb Z$-代数とし，
次数3の正則正規族 $g_i\in\mathcal B_{i,i+3}$ を持つとする。
このとき
\begin{equation}
 \mathcal B\simeq\mathcal B(1).
\end{equation}
従って通常の connected graded 5次元AS正則代数 $D$ により
$\mathcal B\simeq\check D$ と書ける。
特に $\dim\mathcal B_{i,i+n}=\binom{n+4}{4}$ ならば
$H_D(t)=(1-t)^{-5}$ である。
\end{corollary}
\begin{proof}
定理\ref{thm:moduli-ambient-gcd}を $N=5,d=3$ に適用する。
実際，\eqref{eq:moduli-ambient-sigma}の $\sigma$ と
\eqref{eq:moduli-ambient-rho}の $\rho$ から
\begin{equation}
 \theta_{ij}=
 \rho^{-1}_{i+1,j+1}\circ
 \sigma_{i+3,j+3}\circ\sigma_{ij}:
 \mathcal B_{ij}\longrightarrow\mathcal B_{i+1,j+1}
\end{equation}
が1周期同型となる（$2\cdot3-5=1$）。
$\mathcal B_{0n}$ 上の積を
$a*b=a\,\theta^m(b)$（$a\in\mathcal B_{0m}$）と定めれば，
所望の次数付き代数が復元される。
\end{proof}

\begin{proposition}[低次数の核からの正規性判定]\label{prop:moduli-cubic-normal-criterion}
$\mathcal B$ は次数一生成で
$\dim\mathcal B_{i,i+1}=5$ とする。
斉次両側イデアル $I\subset\mathcal B$ が
\begin{equation}
 I_{i,i+3}=\CC g_i,\qquad \dim I_{i,i+4}=5
\end{equation}
を満たし，各 $g_i$ が左右正則なら，$(g_i)$ は次数3の正則正規族である。
\end{proposition}
\begin{proof}
$I_{i,i+4}$ の二つの部分空間
\begin{equation}
 g_i\mathcal B_{i+3,i+4},
 \qquad \mathcal B_{i,i+1}g_{i+1}
\end{equation}
は，正則性によってともに5次元であり，よって一致する。
これは長さ一の正規関係である。
次数一生成と結合律による帰納法で任意の $j\geq i$ に対する
\eqref{eq:moduli-ambient-normal-family}を得る。
\end{proof}

標準多項式型の行Hilbert級数$(1-t)^{-5}$と
次数3の正則元による商の級数$(1-t^3)/(1-t)^5$が既知なら、
次数3と4の差から
$\dim I_{i,i+3}=35-34=1$、$\dim I_{i,i+4}=70-65=5$を計算できる。
ただしHilbert級数だけでAS正則性や正則性を導いたわけではない。
一般に$N=n+2,d=3$なら$n=2,3$では1周期性が従う一方、
$n=4$では最小周期が1または3を割るとしか言えず、3周期の非1周期例の存在は主張しない。

\begin{remark}[適用範囲と既知結果との関係]\label{rem:moduli-ambient-pair-scope}
本稿で構成した中心三次対の $B_F$ と
$G_F\in Z(B_F)_3$ は，初めから通常の
次数付き代数と中心正則元なので，その付随 $\mathbb Z$-代数は
明らかに1周期的である。
系\ref{cor:moduli-cubic-ambient-periodic}の意義は，
\emph{正則正規三次超曲面表示だけを仮定し，
ambientに1周期構造を選ばなくても同じ結論が従う}点にある。

一方，構成した $\theta$ が商イデアル $I=(g_i)$ を保ち，
$\theta(I)=I(1)$ を満たすとはこの議論だけでは言えない。
従って，商 $\mathcal B/I$ が同時に一つの次数付き商 $D/(G)$ として
表示されることや，$G$ を中心化できること，$D$ が
Calabi--Yauとなることは別の問題である。
抽象的な非可換cubicの全てがこの超曲面表示を持つことも主張しない。
この周期性定理は既知のKoszul--Frobenius双対性と正規族の性質の
組合せであり，新規性・初出の最終確認は未了である。
\end{remark}


\section{二次包絡によるモジュライ同型類の判定}\label{sec:moduli-iso}
\begin{theorem}[三次商からのambient復元]\label{thm:moduli-iso}
$B=T(V)/(R_2)$を次数一生成の二次連結次数付き代数とし、
$G\in Z(B)_3$を非零正則中心元とする。$A=B/(G)$に対して
\begin{equation}\label{eq:moduli-envelope}
 B\simeq T(A_1)/
 \left\langle\ker(A_1\otimes A_1\longrightarrow A_2)\right\rangle,
 \qquad \ker(B_3\longrightarrow A_3)=kG
\end{equation}
となる。特に第\ref{sec:normal}節の17係数一般族では、
中心三次対の同型類の次元17が、実際の$A$の
次数付き同型類の次元でもある。
さらに$A$が左右noetherianなAS Gorenstein代数で
injective dimensionが2以上なら、
$(\qgr A,\OO_A,(1))$から同じ対象を復元できる。
\end{theorem}
\begin{proof}
$G$は次数三なので$A_1=B_1$, $A_2=B_2$である。
したがって$A$の二次関係の核は$R_2$であり、
二次関係のみで定義された$B$を復元する。
商写像の次数三の核は$B_0G=kG$である。
次数付き同型$A\simeq A'$はこの復元により
$(B,kG)\simeq(B',kG')$を誘導する。
偏極付き圏については、AS Gorenstein性より
$\operatorname{Ext}_A^0(k,A)=\operatorname{Ext}_A^1(k,A)=0$であり、
Artin--Zhangのsection algebraの自然写像は
\begin{equation}\label{eq:moduli-section}
 A\xrightarrow{\ \sim\ }\bigoplus_{n\ge0}
 \operatorname{Hom}_{\qgr A}(\OO_A,\OO_A(n))
\end{equation}
を与える。よって偏極付き圏の同値も次数付き代数の
同型を誘導する。
\end{proof}

% Source recovery: noncommutative central cubic pair, all Artin bases.
% Main mathematical input: original 2026-11-01 paper V, sections 12--13 and 21.
% This section is included in papers/nc_cubic_threefold_moduli.tex.
\section{中心三次対の全Artin形式局所構造と二次包絡}
\label{sec:moduli-artin}
本節は、滑らかな可換三次三様体の近傍において、
有限係数の17次元族と異なる\emph{完備形式局所モジュライ}を扱う。
以前の中心三次対原稿の第12--13節から定理・証明の有効部分を
再構成したものである。主張は生成空間に標識を入れた
通常の集合値Artin変形関手に関するものであり、
derived stackの同型でも非可換領域全体の大域分類でもない。
外部の同次formality、Ore局所化およびAbel圏変形の
適用条件について、原稿に記載した証明と今回の独立監査を区別する。

以下$k=\C$、$V=k^5$、$S=\Sym V^*$、
$F\in S_3$を滑らかな三次三様体
$X=V_+(F)\subset\PP(V)$の定義式とする。
$(R,\mathfrak m)$は剰余体$k$の局所Artin代数である。

\begin{lemma}[平坦な正規三次商は自動的に中心商]
\label{lem:artin-normal-central}
$B_R$を$R$上各次数有限自由な連結非負次数付き代数とし、
一次で生成され、特殊ファイバー$B_R/\mathfrak mB_R=S$とする。
$w\in(B_R)_3$が$F$に特殊化するなら$w$は左右正則であり、
以下の条件は同値である：
\[
 B_R/(w)\text{ が各次数で }R\text{-平坦},\qquad
 wB_R=B_Rw,\qquad
 w\in Z(B_R).
\]
成立するとき商のHilbert級数は$(1-t^3)/(1-t)^5$である。
\end{lemma}
\begin{proof}
$w$の左右乗法の各次数成分は、$\mathfrak m$で還元すると
$F:S_{d-3}\hookrightarrow S_d$である。
$F\ne0$と$S$が整域であることから単射となり、
有限自由行列の最大rank小行列式が単元であるため、
元の乗法も分裂単射になる。

$J=(w)$について商が平坦なら
$0\to J_d\to(B_R)_d\to(B_R/J)_d\to0$は
有限自由$R$加群の分裂完全列である。
$J_d/\mathfrak mJ_d=FS_{d-3}$なので、
左右乗法$(B_R)_{d-3}\to J_d$はNakayamaと
先の単射性により同型である。従って$J=wB_R=B_Rw$。
正規性から$wb=\sigma(b)w$で次数付き自己同型$\sigma$が得られる。
特殊ファイバーでは$\bar\sigma=\mathrm{id}$であり、
$b=w$と右消去によって$\sigma(w)=w$となる。
$\sigma\equiv\mathrm{id}\pmod{\mathfrak m^r}$とし、
$\mathfrak m^r/\mathfrak m^{r+1}=I$とおくと、
次の差分は次数零の導分
\[
 D:S\longrightarrow I\otimes_k S,\qquad
 D(F)=0
\]
である。$D=\sum_i\ell_i\partial_{x_i}$、
$\ell_i\in I\otimes S_1$だから
$\sum_i\ell_iF_i=0$。滑らかさから
$F_0,\ldots,F_4$は二次正則列をなし、
Koszul複体に一次係数のsyzygyが存在しない。
よって$D=0$となり、$\mathfrak m$の冪零性から
$\sigma=\mathrm{id}$を得る。
中心性から商平坦性へは、
各次数の分裂単射$w:(B_R)_{d-3}\to(B_R)_d$を使えばよい。
\end{proof}

\begin{corollary}[標識付き対の自己同型]
\label{cor:artin-scalar}
上の中心正規対$(B_R,Rw)$の次数付き自己同型が
特殊ファイバーで恒等なら、次数一の生成元への作用は
すべて同じ単元$\lambda\in1+\mathfrak m$によるスカラー倍である。
従って射影標識を入れた対には非自明な無限小自己同型はない。
\end{corollary}
\begin{proof}
$\phi(w)=cw$、$c\in1+\mathfrak m$と書き、
$c$の一意の三乗根による次数スカラー倍を合成すると
$w$を固定する自己同型に帰着する。
その$\mathfrak m$進差分は前補題の$D(F)=0$と同じ
一次Koszul syzygyで消える。
\end{proof}

\subsection{相対Casimir条件と五つのPl\"ucker方程式}
体積形式$\Omega=dx_0\wedge\cdots\wedge dx_4$を固定する。
$F_R\in R\otimes S_3$を$F$の持上げとし、
$a=\sum_{i<j}a_{ij}\,dx_i\wedge dx_j$を
係数$a_{ij}\in\mathfrak m$の定数二形式とする。
\[
 \iota_{\Pi(F_R,a)}\Omega=dF_R\wedge a
\]
で二次bivectorを定める。
$\{i,j,k,l\}=\{0,\ldots,4\}\setminus\{r\}$、
$i<j<k<l$に対し
\[
 q_r(a)=(-1)^r
 (a_{ij}a_{kl}-a_{ik}a_{jl}+a_{il}a_{jk}),\qquad
 g=\sum_{r=0}^4q_r(a)\,\partial_rF_R
\]
とおく。

\begin{lemma}[全Artin上でのCasimir核]
\label{lem:artin-casimir}
二次bivector$\Pi$が$\iota_{dF_R}\Pi=0$を満たすことと、
一意な$a\in R\otimes\wedge^2V^*$を用いて
$\Pi=\Pi(F_R,a)$と表せることは同値である。
この同型は任意の$R$の基底変換と可換であり、
$\Pi$は発散零である。
\end{lemma}
\begin{proof}
$F$の五つの偏微分は二次正則列だから、そのKoszul複体の
\[
 \wedge^3V\otimes R
 \xrightarrow{\iota_{dF_R}}
 S_{R,2}\otimes\wedge^2V
 \xrightarrow{\iota_{dF_R}}
 S_{R,4}\otimes V
\]
の該当次数は、特殊ファイバーで左端の単射と中央での完全性を持つ。
有限自由加群の複体なので、分裂小行列式の単元性を
Artin基底へ持ち上げて同じ完全性を得る。
$\wedge^3V\simeq\wedge^2V^*$により表示と一意性が従う。
$d(dF_R\wedge a)=0$から発散零である。
\end{proof}

\begin{proposition}[JacobiイデアルとPl\"uckerイデアルの一致]
\label{prop:artin-pfaffian}
$\Pi=\Pi(F_R,a)$に対し、Jacobi方程式の係数イデアルは
\[
 (q_0(a),\ldots,q_4(a))
\]
と一致する。特に集合論的なrank分類だけでなく、
非被約な局所Artin基底上での方程式を同定する。
\end{proposition}
\begin{proof}
$J=\frac12[\Pi,\Pi]$と書くと、五変数の外積計算から
\[
 \iota_{\Pi\wedge\Pi}\Omega=2g\,dF_R,\qquad
 \iota_J\Omega=dg\wedge dF_R
\]
が得られる。これは多項式恒等式である。
符号は$ a=dx_0\wedge dx_1+dx_2\wedge dx_3$
の代入で直接確認できる。
線形写像
\[
 L_F:k^5\longrightarrow\Omega^2_{S,3},\qquad
 (c_r)\longmapsto d\bigl(\textstyle\sum_rc_rF_r\bigr)\wedge dF
\]
は単射である。核の元について、
$F_i$のKoszul完全性から
$d(\sum c_rF_r)$は$dF$の多項式倍となるが、
係数の次数は負なので零である。
Euler式と定数係数のKoszul syzygyの消滅により$c_r=0$。
従って$L_F$の非零$5$次小行列式は
$R$上でも単元であり、$L_{F_R}$は分裂単射となる。
上の$J$の表示の係数から$q_r(a)$を回収でき、
両方のイデアルが一致する。
\end{proof}

\begin{theorem}[中心三次対の全Artin形式局所環：入力を明示した結論]
\label{thm:artin-central-pair}
局所Artin基底上の次数別平坦な二次多項式環の標識付き変形と
同次Poisson変形との間に、固定したGL$_5$同変の
Kontsevich形式量子化を適用する。
中心三次直線が$kF$へ特殊化する対の古典的変形関手は
\[
 \widehat{\mathcal I}_{F}\simeq
 \operatorname{Spf}\frac{
  k[[s_1,\ldots,s_{34},a_{01},\ldots,a_{34}]]
 }{(q_0(a),\ldots,q_4(a))}
\]
で表される。形式局所次元は$41$、接空間次元は$44$。
特殊ファイバーで恒等に還元する射影座標変更を割ると、
\[
 \operatorname{Spf}\frac{
 k[[t_1,\ldots,t_{10},a_{01},\ldots,a_{34}]]
 }{(q_0(a),\ldots,q_4(a))}
\]
となり、次元$17$、接空間次元$20$である。
有限群$\Aut(X)$はこの記述ではまだ割らない。
\end{theorem}
\begin{proof}
同次形式量子化の星積を$B_\Pi$とする。
$\Pi$に対するtwisted Taylor mapの次数三部分
$D^0_\Pi:S_{R,3}\to(B_\Pi)_3$は
還元が恒等の有限自由加群の同型である。
また次数四のベクトル場$Y$の評価
\[
 D^1_\Pi(Y)\longmapsto
 \bigl(D^1_\Pi(Y)(x_0),\ldots,D^1_\Pi(Y)(x_4)\bigr)
\]
は同じく還元が恒等の同型なので、
twisted chain identityにより
\[
 D^0_\Pi(F_R)\in Z(B_\Pi)
 \quad\Longleftrightarrow\quad [\Pi,F_R]=0
\]
となる。ここで中心の形成と非平坦基底変換の交換性は
仮定せず、有限自由行列の可逆性によって逆向きを証明する。
このformalityの恒等式は
\href{https://arxiv.org/abs/math/0012228}{Cattaneo--Felder--Tomassini}
のtwisted formulaを入力としている。
補題\ref{lem:artin-normal-central}と
命題\ref{prop:artin-pfaffian}により、
$F_R$の射影変形の34変数と$a$の10変数、
五つのPl\"ucker方程式が得られる。
後者は$\mathrm{Gr}(2,5)$の非零アフィン錐の閉包で
次元$7$、原点の接空間次元$10$なので
$34+7=41$、$34+10=44$。
$X$が滑らかで$H^0(X,T_X)=0$だから
$\PGL_5$の軌道方向24は形式的に横断sliceから分離できる。
$34-24=10$を得て後半が従う。
なお特異な形式基底でのCalabi--Yau性は相対的であり、
絶対大域次元が五であるとは主張しない。
\end{proof}

\begin{remark}[一次方向はすべて二次へ延長するわけではない]
\label{rem:artin-obstruction}
$a=\varepsilon(dx_0\wedge dx_1+dx_2\wedge dx_3)$を
$k[\varepsilon]/(\varepsilon^2)$上で考えると
$q_r(a)=0$だが、
$k[\varepsilon]/(\varepsilon^3)$上では$q_4(a)=\varepsilon^2$。
二次の補正$b$を$a+\varepsilon^2b$に加えてもこの項は消えない。
従って接空間20次元と局所次元17の差は実際の二次障害を表し、
全Artin比較を無障害性と混同してはいけない。
\end{remark}

\subsection{Ore局所化、偏極、および斉次二次包絡}
次数別平坦な斉次変形$A_R$を
$S/(F)$へ特殊化するものとする。
一次元の持上げ$\widetilde x_i$を選び、
標準開集合$U_i=D_+(x_i)\cap X$上で
\[
 \mathscr A_R(U_i)=
 \bigl(A_R[\widetilde x_i^{-1}]\bigr)_0
\]
を取る。$[A_R,A_R]\subset\mathfrak m A_R$かつ
$\mathfrak m$冪零なので、有限の交換子展開により
$\{\widetilde x_i^n\}$は左右Ore集合となる。
その元の特殊化は$S/(F)$の整域の非零因子だから、
各局所化は$R$平坦である。
重なりは二つの一次元をともに可逆化した次数零部分で
記述され、三重重なりのcocycleも局所化の普遍性に従う。
この操作で$X$上に平坦な結合$R$代数層
$\mathscr A_R$が得られる。

\begin{proposition}[偏極とsection algebraの回復]
\label{prop:artin-ore-section}
上のOre局所化層の次数$n$部分を貼り合わせた
$\mathscr E_n$は可逆両側$\mathscr A_R$加群であり、
\[
 \mathscr E_m\otimes_{\mathscr A_R}\mathscr E_n
 \simeq\mathscr E_{m+n}.
\]
$n\ge0$について
\[
 H^i(X,\mathscr E_n)=0\ (i>0),\qquad
 (A_R)_n\xrightarrow{\sim}H^0(X,\mathscr E_n)
\]
が成り立ち、積とも両立し任意のArtin基底変換に可換である。
\end{proposition}
\begin{proof}
局所化の一次元は可逆なので強い$\mathbb Z$次数付けと
両側可逆性が従う。有限標準被覆の\v{C}ech複体を取り、
$R$の小拡大を順に帰納すると、
特殊ファイバーの$H^i(X,\OO_X(n))=0$ ($i>0,n\ge0$)
から高次消滅と$H^0$の有限自由性が得られる。
特殊化で同型な$(A_R)_n\to H^0(X,\mathscr E_n)$は
Nakayamaで同型となる。強次数付けにより積との両立も従う。
\end{proof}

\begin{lemma}[可逆両側偏極の同型類の一意性]
\label{lem:artin-polarization}
$\mathscr A_R$上の二つの平坦な可逆両側加群
$\mathscr M,\mathscr N$が同じ線束$L$に特殊化し、
$H^1(X,\OO_X)=H^0(X,T_X)=0$なら、
$\mathscr M\simeq\mathscr N$である。
$H^0(X,\OO_X)=k$なら各々の自己同型は$R^\times$である。
\end{lemma}
\begin{proof}
小拡大$R'\twoheadrightarrow R$の核$I$で
$\mathfrak mI=0$とする。
特殊ファイバーの局所基底の持上げにより
$\mathscr M,\mathscr N$は右加群として局所自由階数一。
同型の右加群としての持上げの障害は
$I\otimes H^1(X,\OO_X)$に属し消える。
二つの左作用の差は導分に対応し
$I\otimes H^0(X,T_X)=0$なので持上げは両側同型。
自己同型群は
$H^0(X,\OO_X)=k$と小拡大帰納から$R^\times$となる。
\end{proof}

\begin{theorem}[斉次環の全Artin変形の二次cover：比較定理を入力]
\label{thm:artin-quadratic-cover}
上の中心三次対の形式関手と、
$X$上の平坦結合代数層の標識付き古典的Artin変形関手の間の
Ore局所化による自然変換$\Theta$を考える。
YekutieliのArtin変形量子化の自然性と
De~Deken--Lowenの$\mathbb Z$-生成列の変形比較を
この設定に適用すると、$\Theta$は自然同型となる。
従って任意の次数別平坦な斉次変形$A_R$は
\[
 V_R=(A_R)_1,\quad
 K_{2,R}=\ker(V_R^{\otimes2}\to(A_R)_2),\quad
 B_R^{\mathrm{quad}}=T_R(V_R)/(K_{2,R})
\]
を用いて一意の中心三次対から回復される。
特に$B_R^{\mathrm{quad}}$は次数別$R$平坦で、
\[
 L_R=\ker\bigl((B_R^{\mathrm{quad}})_3\to(A_R)_3\bigr)
 \subset Z(B_R^{\mathrm{quad}})_3
\]
は階数一の直和因子であり、
$A_R\simeq B_R^{\mathrm{quad}}/(L_R)$。
この構成はArtin基底変換と可換である。
\end{theorem}
\begin{proof}[原研究稿の証明構造]
まず局所化による$\Theta$を上で実際に構成した。
中心対の形式関手は
定理\ref{thm:artin-central-pair}の完備局所環で表される。
構造代数層のArtin変形関手も、
\href{https://arxiv.org/abs/0905.0488}{Yekutieli, Cor.~11.2}
の量子化と可換三次の変形sliceから同じ環で表される。
接写像$d\Theta$は、可換方向では
$H^1(X,T_X)$上の恒等写像であり、
反対称交換子方向では
\[
 \wedge^2V^*\longrightarrow H^0(X,\wedge^2T_X),
 \qquad a\longmapsto\iota_Ea
\]
という10次元空間の同型である。
よって接写像は可逆な三角行列となる。
完備局所環の連続自己写像が余接空間で同型なら
各$\mathfrak m^r/\mathfrak m^{r+1}$上でも同型であり、
形式的Nakayama帰納から環自体が同型となる。
したがって$\Theta$は全Artin基底で自然同型。

一方、上の命題で$\mathscr A_R$の局所化次数一は
$\OO_X(1)$の両側可逆な持上げ$\mathscr E_1$となり、
全section algebraを回復する。
補題\ref{lem:artin-polarization}でこの持上げの同型類は
一意なので、任意の$A_R$と$\Theta$から来る中心商を
次数付き環として同一視できる。
中心商では三次まで新しい二次関係がないから
二次包絡は元の$B_R$に一致し、三次核は$L_R$。
この方法は未知の二次包絡の平坦性を先取りせず、
既に存在する中心対に同型化した後でそれを導く。
基底変換は有限自由の分裂核と局所化の普遍性で従う。

ここで使用した一次資料は
\href{https://arxiv.org/abs/1209.6243}{Yekutieli,
 \emph{Deformations of affine varieties and the Deligne crossed groupoid}}
のOre局所化、
\href{https://arxiv.org/abs/math/0511197}{Lowen,
 \emph{Algebroid prestacks and deformations of ringed spaces}}、
\href{https://arxiv.org/abs/1001.4265}{De~Deken--Lowen,
 \emph{Abelian and derived deformations in the presence of
 $\mathbb Z$-generating geometric helices}}である。
これらの定理の適用範囲の独立した一次資料監査は未完成である。
\end{proof}

\begin{remark}[明確に主張しない事柄]
\label{rem:artin-limits}
この全Artin同型は偏極付き\emph{古典的}変形関手の同定であり、
derived moduli stack、高次自己同型を保つgroupoid、
構造対象を忘れたKuznetsov成分の変形、
またはそこで構成される可換復元との一致を与えない。
無限小の中心対を得ることと、有限非可換パラメータで
滑らかなtailsを持つことは別問題である。
\end{remark}

\subsection{固定ambientからの忘却の四次元ファイバー}
\begin{proposition}[中心三次直線の射影束]
\label{prop:artin-fixed-ambient-fiber}
ある良い開集合$U$の各点$B$が二次Koszul AS正則CY五次元整域であり、
$\dim Z(B)_1=2$、$\dim Z(B)_3=5$かつ
中心の形成が基底変更と可換であると仮定する。
$E_1=Z(\mathscr B)_1$、$E_3=Z(\mathscr B)_3$とおけば
\[
 0\longrightarrow\Sym^3E_1\longrightarrow E_3
 \longrightarrow L\longrightarrow0,\qquad\rank L=1.
\]
中心三次直線のパラメータ空間は
$\PP_U(E_3)\to U$で各fiberは$\PP^4$。
$\PP(E_3)\setminus\PP(\Sym^3E_1)$は
$\Sym^3E_1\otimes L^{-1}$を作用群とするtorsorである。
同型類を取る場合、$\Gamma_B=\Aut_{\gr}(B)/k^\times$
が有限なら一般fiberの粗い双有理型は
$\PP Z(B)_3/\Gamma_B$となり、その次元は4である。
\end{proposition}
\begin{proof}
中心三次部のベクトル束の射影化で前半を得る。
$B$が整域なので中心三次の非零元は正則、
次数ごとの乗法の核は零で、商の各次数が局所自由となる。
$E_3\to L$に非零に写る直線はこの短完全列のsplittingに対応し、
splittingの差は$\Hom(L,\Sym^3E_1)$に属する。
同型類については、二次包絡と三次核の復元から
$B/(W)\simeq_{\gr}B/(W')$と
$[W'],[W]$が$\Gamma_B$軌道に属することが同値である。
Rosenlichtの有理商を$U$とincidenceへ適用し、
周囲$B$の軌道の上の空間が
$\PGL(V)\times_{\Gamma_B}\PP Z(B)_3$となることから
不変関数体が$k(\PP Z(B)_3)^{\Gamma_B}$に一致する。
有限群商は次元を保つ。
\end{proof}

\begin{remark}[四係数は標識を忘れても消えない]
特に固定した周囲に対し$Z(B)_3=kG\oplus\Sym^3Z(B)_1$なら、
$G+f(u,v)$の四係数が相対アフィンfiberを表す。
有限群$\Gamma_B$による同一視は生じ得るが、
四係数全体を座標変換で消せるわけではない。
原稿中の明示剛体周囲$B_*$についての
$\Aut_{\gr}(B_*)=k^\times$と正則開集合の主張は、
付属の有限整数証明書が現行Git treeに存在しないため、
本稿では独立に再検証した定理として採用しない。
\end{remark}

% Part of the present nc cubic threefold moduli manuscript.
% Restored from the 2026-11-01 central-pair paper V, sections 19--21.
\section{全混合族の低次数中心、偶数Veroneseと剛体ambient}
\label{sec:moduli-rigid-ambient}
本節では、可換点の全Artin比較とは別に、有限非可換パラメータ
において一つのambientを固定したまま三次方程式を動かす問題を調べる。
元の中心三次対原稿にある命題を現在有効な仮定と証明に整理する。
非可換Bertiniの適用、三次Sklyaninの位数条件、
および有限整数証明書の独立監査は未完了である。

\subsection{三次商による二次中心の非増大}
$U=\Span(u,v)$、$R=k[u,v]$、$k=\C$とおく。
$F\in k[x,y,z,u,v]_3$は
$F(x,y,z,0,0)$が滑らかな平面三次式$f_3(x,y,z)$であると仮定する。
$h\in k^\times$に対して
\[
 B(h,F)=\frac{R\langle x,y,z\rangle}
 {([x,y]-h\Symm F_z,\ [y,z]-h\Symm F_x,\
    [z,x]-h\Symm F_y)}
\]
とし、次数三の中心正則元$W$で割った
$A=B/(W)$を考える。ここで$B$は一次関係を持たないとする。

\begin{lemma}[中心三次商では二次中心が増えない]
\label{lem:moduli-no-new-quadratic-center}
上の仮定で、自然な同一視$B_2\simeq A_2$の下で
$Z(A)_2=Z(B)_2$となる。
\end{lemma}
\begin{proof}
$Q\in B_2$の像が$A$で中心なら、
$[Q,x_i]=c_iW$、$c_i\in k$、
$(x_1,x_2,x_3)=(x,y,z)$と書ける。
$u,v,W$が中心なので
\[
 0=[Q,[x_i,x_j]]
 =h[Q,\Symm F_{x_k}]
 =hW\sum_\ell c_\ell F_{x_kx_\ell}
\]
である（$i,j,k$は巡回）。
$hW$を消去すると、右側の一次式が$B_1$で零。
一次関係を持たないから可換多項式として零であり、
$u=v=0$で
\[
  \sum_\ell c_\ell(f_3)_{x_kx_\ell}=0
  \quad(k=1,2,3).
\]
すなわち$\sum_\ell c_\ell(f_3)_{x_\ell}$は
定数の斉次二次式なので零になる。
$c\ne0$なら座標を取り換えれば$f_3$は一変数を含まず、
対応する射影点で偏微分がすべて消えて滑らかさに矛盾する。
従って$c=0$で、$Q$は既に$B$の中心である。
\end{proof}

\begin{lemma}[全混合族の次数一・二の中心]
\label{lem:moduli-low-center}
$S=B/(u,v)$とし、中心三次元$\gamma\in S_3$で
\[
 S/(\gamma)\simeq B(E,\mathcal L,\tau)
\]
となるとする。$E$は滑らかな楕円曲線、
$\deg\mathcal L=3$、$\tau$は平行移動で$\tau^2\ne1$と仮定する。
さらに相対PBWにより
$\gr_{x,y,z} B\simeq S\otimes_k k[u,v]$とする。
このとき
\[
 Z(B)_1=Z(A)_1=U,\qquad
 Z(B)_2=Z(A)_2=\Sym^2U.
\]
\end{lemma}
\begin{proof}
twisted homogeneous coordinate ringは
$k(E)[t,t^{-1};\tau]$に埋め込まれる。
そこに次数$m=1,2$の中心元$at^m$があれば
一次切断$s_1,s_2$との可換性から
$\tau^m(s_1/s_2)=s_1/s_2$を得る。
次数三の完全線形系の一次切断の比は$k(E)$を生成するため
$\tau^m=1$となり仮定に反する。
従って$S/(\gamma)$には低次数中心がなく、
$S\to S/(\gamma)$は次数一・二で単射だから
$Z(S)_1=Z(S)_2=0$。
$B$の中心元の最高相対次数を取ると、
正の相対次数一・二は消え、中心は$k[u,v]$に属する。
三次商でも次数一の関係は増えず、
二次は補題\ref{lem:moduli-no-new-quadratic-center}を使う。
\end{proof}

\subsection{偶数Veroneseからの標識付き代数の回復}
\begin{lemma}[中心二変数の平方根と奇数部分]
\label{lem:moduli-even-veronese}
$A$はconnectedな次数付き$k$代数で、
二つの一次中心元$u,v$が多項式環$S=k[u,v]$を生成し、
$A$が$S$上自由、
$Z(A)_2=S_2$とする。
$R=A^{(2)}=\bigoplus_{n\ge0}A_{2n}$とおくと、
$R$の次数付き同型類は$A$の次数付き同型類を決定する。
詳しくは、同じ仮定を満たす$A'$について、
任意の$\phi:R\xrightarrow{\sim}(A')^{(2)}$は
正確に二つの次数付き同型$A\to A'$へ持ち上がり、
その差は奇数次数に$-1$を掛ける自己同型である。
\end{lemma}
\begin{proof}
$z\in Z(A)_1$について$uz,vz\in S_2$。
$u(vz)=v(uz)$と$S$の一意分解性から
$uz=u\ell$、$\ell\in S_1$であり、
$u$の正則性から$z=\ell$となる。
$z\in Z(R)$と奇数次数の$a\in A$に対しては
$ua\in R$なので
$u[z,a]=[z,ua]=0$より$[z,a]=0$。
従って$Z(R)=Z(A)^{(2)}$。
その次数一中心の生成する内在的部分環は
\[
 D=k[u^2,uv,v^2]
   \simeq k[\alpha,\beta,\gamma]/(\alpha\gamma-\beta^2).
\]
$\Proj D$は二次曲線で$\OO_{\Proj D}(1)=\OO_{\PP^1}(2)$。
任意の$D\simeq D'$は二次Veronese表現
$\GL_2\to\Aut_{\gr}D$により二通りだけ
$S\simeq S'$へ持ち上がり、その差は全一次元の符号反転。

$u,v$を局所化して$R$の次数を二倍すると
\[
 A[u^{-1}]
  \simeq R[(u^2)^{-1}]
    \otimes_{D[(u^2)^{-1}]}S[u^{-1}]
\]
であり、$v$の場合も同様である。
$A$は$S$上自由なので
\[
 A[u^{-1}]\cap A[v^{-1}]=A
 \quad\text{in }A[(uv)^{-1}].
\]
これが$R$と$S$だけからの復元式である。
同型を持ち上げた後、偶数次数と$S_1$上で恒等なら、
奇数元$a$について$u\phi(a)=\phi(ua)=ua$から
$\phi(a)=a$。残る一つは全奇数次数の符号反転である。
\end{proof}

\begin{theorem}[構造対象を保つ導来同値：条件付き全混合復元]
\label{thm:moduli-marked-derived}
$A=B/(W)$と$A'=B'/(W')$は上の全混合族の中心正則三次商で、
補題\ref{lem:moduli-low-center}の仮定と
それぞれの$k[u,v]$上の自由性を満たすとする。
両方の$\qgr$は有限大域次元を持つと仮定する。
次は同値である：
\begin{enumerate}
\item $A\simeq_{\gr}A'$;
\item $\OO_A$を$\OO_{A'}$に送る$k$線形Abel圏同値
 $\qgr A\simeq\qgr A'$が存在する;
\item ある$m,r\in\mathbb Z$について
 $\OO_A$を$\OO_{A'}(m)[r]$に送る$k$線形三角同値
 $D^b(\qgr A)\simeq D^b(\qgr A')$が存在する。
\end{enumerate}
これらの条件は、周囲の二次代数と三次中心直線
$(B,kW)$の同型類も決定する。
\end{theorem}
\begin{proof}
$(1)\Rightarrow(2)\Rightarrow(3)$は標準的である。
$(3)$の同値にshiftとtwistを合成して構造対象を一致させる。
$B$はCY五次元、$W$は中心正則三次元だから
balanced dualizing complexは$A(-2)[4]$である。
tailsの有限大域次元を使ってSerre関手は
$\mathsf S_A\simeq(-2)[3]$となる。
これは\href{https://www.math.bgu.ac.il/~amyekut/publications/nc-serre/paper.pdf}
{Yekutieli--ZhangのSerre双対}の適用を入力とする。
同値はSerre関手と可換なので
$\mathsf H_A=\mathsf S_A^{-1}[3]\simeq(2)$を保ち、
その軌道環
\[
 \bigoplus_{n\ge0}\Hom(\OO_A,\mathsf H_A^n\OO_A)
 \simeq A^{(2)}
\]
を積とともに同定する。
補題\ref{lem:moduli-low-center}と
\ref{lem:moduli-even-veronese}により$A\simeq A'$。
最後に二次関係の核から$B$を、三次の核から$kW$を回復する
定理\ref{thm:moduli-iso}を使う。
\end{proof}
\begin{remark}
\label{rem:moduli-marking}
ここで除いたのは次数一shiftの\emph{指定}であり、
構造対象の条件自体は残る。任意の導来自己同値が
構造対象をline objectのshiftへ送るとは主張しない。
有限大域次元を持たないtailsへの適用も主張しない。
\end{remark}

\subsection{全混合ポテンシャルの有限同型判定}
\begin{proposition}[中心blockと三次ポテンシャルの共変性]
\label{prop:moduli-mixed-block-isomorphism}
上の低次数中心を持つ$B(h,F)$、$B(h',F')$について
任意の次数付き同型$B(h',F')\to B(h,F)$は
\[
 s' \longmapsto Ns,\qquad x'\longmapsto Lx+Ms,
 \quad
 L\in\GL_3(k),\ N\in\GL_2(k),\ M\in\mathrm{Mat}_{3\times2}(k)
\]
の形を持つ。これは同型である必要十分条件が
ある$r\in k[u,v]_3$について
\[
 h'F'(Lx+Ms,Ns)=h\det(L)F(x,s)+r(s)
\]
となることに等しい。
中心商の次数付き同型にはさらに
$\psi_{L,M,N}(W')=\kappa W$ ($\kappa\in k^\times$)
を要し、これで必要十分である。
\end{proposition}
\begin{proof}
低次数中心が$U=\Span(u,v)$なので同型はblock形となる。
$([y,z],[z,x],[x,y])^t$の変換は
$\det(L)L^{-t}$であり、関係式の保存条件は
\[
 h'\nabla_{x'}F'(Lx+Ms,Ns)
 =h\det(L)L^{-t}\nabla_xF(x,s).
\]
二次関係空間が交代二次式から対称二次式へのgraphであるため、
これは可換二次多項式の恒等式と同値。
chain ruleで勾配の式を積分すると、差は$k[u,v]_3$に入る。
最後の中心直線保存は二次包絡と三次核の回復に等しい。
\end{proof}

\subsection{射影自己同型のない明示ambient}
次の例は元の中心三次対原稿第21節にあったものであり、
正則な中心商の開集合の存在は、中心商論文の有限位数境界に関する定理から従う。

\begin{proposition}[明示的な剛体ambient：中心の入力を指定]
\label{prop:moduli-rigid-example}
\[
 F_*=\frac{x^3+y^3+z^3}{3}
     +uxy+vyz+(u+v)zx
\]
に対する$h=1$の代数
\[
 B_*=\frac{k[u,v]\langle x,y,z\rangle}
 {([x,y]-z^2-vy-(u+v)x,\
   [y,z]-x^2-uy-(u+v)z,\
   [z,x]-y^2-ux-vz)}
\]
を考える。$S_*=B_* /(u,v)$の中心三次商が、
平面三次$E_*:x^3+y^3+z^3+xyz=0$に付随する
位数6の平行移動を持つTHCRであることを入力とする。
このとき
\[
 Z(B_*)_1=ku\oplus kv,\qquad
 Z(B_*)_3=kG_*\oplus k[u,v]_3,
 \qquad \Aut_{\gr}(B_*)=k^\times.
\]
ここで
\[
\begin{aligned}
 \Omega_1&=\frac{x^3+y^3+z^3}{3}
       -\frac19\sum_{\sigma\in S_3}\sigma(xyz),\\
 G_*={}&\Omega_1+
       u\Symm(xy)+v\Symm(yz)+(u+v)\Symm(zx)\\
 &+\frac1{18}\bigl(vx^2+(u+v)y^2+uz^2\bigr)\\
 &+\frac1{18}\bigl((-u^2-uv+3v^2)x
       +(3u^2+5uv+3v^2)y+(3u^2-uv-v^2)z\bigr).
\end{aligned}
\]
従って$B_*/(G_*+f)$ ($f\in k[u,v]_3$)は、
一つの周囲の中で相互に次数付き非同型な四次元代数族を与える。
さらに、$f$の非空Zariski開部分ではtailsが正則で大域次元3となる。
\end{proposition}
\begin{proof}
Fermat型の混合中心の普遍公式
\href{https://github.com/uedakazushi/nc-cubic/blob/main/papers/support/papers/sections/fermat_mixed_universal_center.tex}
{（nc-cubic論文A）}へ
$A=u,B=v,C=u+v,D=E=H=g=0,h=1$を代入すれば
上記の$G_*$とその中心性が得られる。
指定されたTHCRでは平行移動の位数が6なので
次数一・二の中心は零、次数三では$k\Omega_1$である。
相対PBWと中心元の三次還元から
$Z(B_*)_1=U$、$Z(B_*)_3=kG_*\oplus k[u,v]_3$となる。

任意の次数付き自己同型は中心一次空間を保つから、
$s=(u,v)^t$、$x=(x,y,z)^t$について
$s\mapsto Ns,\ x\mapsto Mx+\ell(s)$の形である。
命題\ref{prop:moduli-mixed-block-isomorphism}の勾配比較より
\[
 F_*(Mx+\ell(s),Ns)=\det(M)F_*(x,s)+r(s)
\]
となる。最高$x$-次数で
$\sum_i(Mx)_i^3=\det(M)\sum_ix_i^3$。
Fermat三次のHessianは非零定数倍の$xyz$なので、
$M$は座標三直線の合併を保ち、monomialである：
$(Mx)_i=\lambda_i x_{p(i)}$。
$x_{p(i)}^2s$項を比較すると$\ell=0$。
$\lambda_i^3=\det(M)$と
$\det(M)=\operatorname{sgn}(p)\prod_i\lambda_i$から
$p$は偶置換である。
同時スカラー倍で$\lambda_1=1$に規格化すると
$\lambda_i^3=1$、$\lambda_1\lambda_2\lambda_3=1$。
混合係数を
$(r_1,r_2,r_3)=(v,u+v,u)$とすれば
$r_i(Ns)=\lambda_i r_{p(i)}(s)$。
これらの一次式の唯一の関係の係数は
$c=(-1,1,-1)$である。よって
$\lambda_i=\mu c_{p(i)}/c_i$。
$p$が三cycleなら三つの比に$+1$と$-1$が混じり、
$\lambda_i^3=1$に矛盾する。
従って$p=1$、$\lambda_i=1$、$N=1$となり、
自己同型は同時スカラー倍だけ。
最後に次数二包絡・三次核の復元により、
$B_*/(G_*+f)\simeq_{\gr}B_*/(G_*+f')$なら
$[G_*+f]=[G_*+f']$。
$G_*$の係数を1と規格化したので$f=f'$である。
\end{proof}

\begin{corollary}[剛体周囲の中の正則な四係数非同型族]
\label{rem:moduli-rigid-generic-smoothness}
ある非零多項式$P(f)$が存在し、$P(f)\ne0$では
$\gldim\qgr B_*/(G_*+f)=3$である。
同時に可換な$V(F_*+f)\subset\PP^4$も滑らかとなる$f$は非空開部分をなし、
その中の有理係数点はZariski稠密である。
\end{corollary}
\begin{proof}
周囲はnoetherian Koszul CY5、相対associated gradedは$S_1\otimes k[u,v]$である。
第\ref{sec:pair-explicit-centers}節の中心式を代入すると$G_*$を得て、
全ての$f$で商は$k[u,v]$上自由である。
最高次中心元は非零正則、境界は位数6の滑らかな楕円THCRなので、
定理\ref{CSBGthm:global-open}の全仮定を満たす。
そこでは全係数chartの代数が$D_u[b,c,d]$となること、
中心持上げにより可換有限型中心上有限となることを確認してからBertiniを適用する。

可換な滑らかさについては$u=v=0$上に勾配の共通射影零点はない。
$u=1,t=v/u$上のactive勾配は
$x^2+y+(1+t)z$, $y^2+x+tz$, $z^2+(1+t)x+ty$であり、
$x,y,z$全次数を先に比較すると互いに素な先頭項$x^2,y^2,z^2$を持つ
monic Gr\"obner基底である。従って臨界スキームは$k[t]$上rank8の有限自由環を持ち、次元1である。
各点が全三様体の特異点となる条件は$(F_*+f)(P)=\partial_t(F_*+f)(P)=0$で、
$f(1,t)$の値と微分は四係数に対する独立な二条件である。
従って特異点incidenceの次元は高々3である。$v=1$も同様。
全incidenceは$\PP^4\times\A^4$内で閉で、その係数空間への像はproper性により閉である。
ゆえに滑らかな係数は非空開をなし、tails正則開部分とも交わる。
有理点の稠密性から最後の主張が従う。
これは可換三様体との残余圏同値を主張するものではない。
\end{proof}


\section{結論と検証の状態}\label{sec:status}
本稿の41次元incidence、17係数正規形、分離型余次元15と
忘却核7次元の主張には本文に証明がある。
全nc cubicモジュライの既約成分性、有限非可換パラメータ全域の変形・障害理論、
一般次元nc cubic moduliの完全分類は未解決。
正規族を持つAS正則$\mathbb Z$-代数のambient周期性は第\ref{sec:moduli-ambient-periodicity}節で証明するが、
その1周期同型が三次イデアルを保つことや、全nc cubicの超曲面表示は未証明。
規格化された37次元CY周囲の被約芽の成分性と一般中心は第\ref{sec:pair-generic-component}節で扱う。全nc cubicの大域成分性とは区別する。
一次文献の初出調査と既存証明の全面独立監査は未完了。
31次行列式、中心の次数別計算、分離型の検算は
\nolinkurl{computations/independent/verify_generic17_slice.py}、
\nolinkurl{computations/independent/verify_generic17_center.py}、
\nolinkurl{computations/verify_split_codimension.py}で再現する。

\appendix
\section{17係数正規形の正確な行列証明書}\label{app:rank}
係数を $\mathbb Q$ 上で計算する．単項式基底は
\[
 x^iy^jz^ku^\ell v^{3-i-j-k-\ell},\quad
 i,j,k,\ell\ge0,\quad i+j+k+\ell\le3
\]
を辞書式の添字順で並べる．$i+j+k>0$ の31単項式だけを行として取り出す．
19個の群方向は $M_{ij}$（$1\le i,j\le3$）９個，
$N_{ij}$（$1\le i,j\le2$）４個，
$C_{ij}$（$1\le i\le3,1\le j\le2$）６個の順に
\eqref{eq:Lie} の係数微分を並べる．このうち最初の１列を除く．
13個の正規形方向は
\[
 x^3+y^3+z^3,\ xyz,\ vy^2,\ (u+v)z^2,\quad
 x(u^2,uv,v^2),\ y(u^2,uv,v^2),\ z(u^2,uv,v^2)
\]
の順である．得られる $31\times31$ 行列の行列式は
\[
 -5538243375000=-2^3\cdot3^5\cdot5^6\cdot7^2\cdot61^2.
\]
付属の \nolinkurl{computations/independent/verify_generic17_slice.py} は \eqref{eq:F0} に対して
これを厳密整数計算で再現する．中心三次方向
$u^3,u^2v,uv^2,v^3$ を最後に付け足すと，
$35\times35$ 行列式も同じ非零値となる．
したがって binary cubic に入る量子的な正規化補正の詳細を仮定せずに
横断性を結論できる．

\section{有限位数楕円境界の中心商に必要な正則性の証明}
\label{CSBGsec:pair-finite-order-background}
この節は固定した周囲の正則な中心商を得るための証明を収める。
一般中心商の分類や全中心の同定は別リポジトリCayley--Sklyanin論文の対象である。
以下で$B$は標数零の体上の左右noetherian connected graded AS正則代数、
$W\in Z(B)$は正次数の正則斉次元、$A=B/(W)$とする。
$\mathcal P_A=\operatorname{thick}\{\pi A(n):n\in\mathbb Z\}$、
$\mathcal Q_A=\mathrm D^b(\qgr A)/\mathcal P_A$と置く。
有限graded自由行列因子化と次数付き特異圏の比較には
\cite[Theorems~4.7, 5.5]{CSBGCCKM}を用いる。
右加群には反対環を適用し、中心元のnormalizing automorphismは恒等である。
従って有限syzygy後のcoherent対象を二周期因子化で扱える。
以下の補題の引用入力と仮定を保ち、一般の非正則底へ拡張しない。
\begin{lemma}[有限射影因子化の縮約と分母払い]
\label{CSBGlem:gluing-contraction}
$2$が可逆な左右noetherian環$S$と中心正則元$w$を取り、
有限生成射影加群の行列因子化を$d^2=w$と書く。
$S/(w)$上のcokernelが有限射影次元なら、この因子化は縮約可能である。
さらに中心正則元$u$で局所化した後に縮約可能なら、
ある$N$について$u^N\mathrm{id}=dK+Kd$となる。
次数付きの場合も、因子化の次数shiftを含めて同じ主張が成立する。
\end{lemma}
\begin{proof}
$\bar S=S/(w)$上の二周期分解から、cokernelの有限射影次元は
各周期syzygyの射影性を意味する。従って分解は分裂し、
$\bar S$上で$d\bar H+\bar Hd=1$となる奇写像を取れる。
有限射影性からこれを$S$上の$H$へ持ち上げると
$dH+Hd=1+wT$となる。
左辺は$d$と可換なので、$w$の正則性により$[d,T]=0$。
ゆえに$H'=H-\frac12dT$は$dH'+H'd=1$を満たす。
局所化後の$H'$の有限個の成分の分母を$u^N$で払い、
$S\hookrightarrow S[u^{-1}]$の単射性を使えば最後の等式を得る。
\end{proof}

\begin{lemma}[有限大域次元からtwistsによる生成へ]
\label{CSBGlem:global-finite-to-perfect}
$A$は次数1生成で、$\operatorname{qgr}A$が有限大域次元を持つなら、
$\mathrm D^b(\operatorname{qgr}A)=\mathcal P_A$である。
\end{lemma}
\begin{proof}
大域次元を$e$とし、graded自由表示を繰り返して
$0\to K\to L_e\to\cdots\to L_0\to M\to0$を取る。
各$L_i$は$\pi A(n)$の有限直和である。
最後のYoneda拡大は$\operatorname{Ext}^{e+1}(M,K)=0$なので、
$M$は有限複体$[L_e\to\cdots\to L_0]$の直和因子である。
line object自体が射影的であるとは仮定していない。
\end{proof}

\begin{proposition}[中心零化元による十分条件]
\label{CSBGprop:global-annihilators}
$u,v\in Z(B)_1$が$A$上正則列をなし、
$\mathrm D^b(\operatorname{qgr}A/(u,v))$が構造対象のtwistsで厚く生成されるとする。
正次数の斉次中心元$p_1,\ldots,p_r\in Z(A)$が
次数付き特異圏の全対象を零化し、ある$N$について
$(u,v)^N\subset(p_1,\ldots,p_r)A$なら、$\mathcal Q_A=0$である。
\end{proposition}
\begin{proof}
graded lift $M$のKoszul複体$K(p_1,\ldots,p_r;M)$のcohomologyは
$(p_1,\ldots,p_r)$に殺され、従って$(u,v)$の冪に殺される。
有限filtrationの各商は$A/(u,v)$加群である。
正則列のKoszul分解から$\pi(A/(u,v))(n)\in\mathcal P_A$なので、
境界の生成性によりこのKoszul複体は$\mathcal P_A$に入る。
一方、$\mathrm D^b(\operatorname{gr}A)\to\mathcal Q_A$は
次数付きperfect複体を殺すから次数付き特異圏を経由する。
従って各$p_i$は$\mathcal Q_A$で零作用し、反復coneは
$M$のshiftとtwistの直和へ分裂する。ゆえに$M$自身も零となる。
\end{proof}

\begin{proposition}[二つの中心chartの貼合せ]
\label{CSBGprop:global-two-charts}
$u,v\in Z(B)_1$が$B$上正則で$A$上正則列をなし、
$A[u^{-1}]_0,A[v^{-1}]_0$が有限大域次元を持つとする。
$\operatorname{qgr}A/(u,v)$の有界導来圏が構造対象のtwistsで厚く生成されるなら、
$\mathcal Q_A=0$である。
\end{proposition}
\begin{proof}
各coherent対象を行列因子化で表し、各chart上で縮約する。
補題\ref{CSBGlem:gluing-contraction}の分母払いから、対象ごとに選んだ
$u^r,v^s$が$\mathcal Q_A$で零作用する。
二重Koszul coneのcohomologyは$(u^r,v^s)$に殺され、
$(u,v)^{r+s-1}\subset(u^r,v^s)$から境界の有限filtrationを持つ。
前命題と同じ議論でconeはperfectである。
中心作用は自然変換なので、最初の零coneの直和分解とも可換し、
二重coneは元の対象の四つのshift・twistの直和となる。
従って元の対象は零である。一様な$r,s$は要求しない。
\end{proof}

\begin{lemma}[perfect性とSerre双対性からの正確な次元]
\label{CSBGlem:gluing-serre-dimension}
$A$は左右noetherian connected graded AS--Gorensteinで、
balanced dualizing complexが$A_\nu(-\ell)[d+1]$（$d\ge0$）で、
$\pi A\ne0$とする。
$\mathrm D^b(\operatorname{qgr}A)=\mathcal P_A$なら
$\operatorname{gldim}\operatorname{qgr}A=d$である。
\end{lemma}
\begin{proof}
\cite[Theorem~4.2, Corollary~4.3]{CSBGYZSerre}は左右の$\chi$と
有限cohomological dimensionを与える。
特に$\pi A(n)$は有限射影次元を持つ。
Chen\cite[Theorem~3.10]{CSBGChenGeneralizedSerre}の一般化Serre双対性は
有限射影次元の対象に適用でき、その関手は
$S(M)=M_\nu(-\ell)[d]$となる。
仮定により全coherent対象が適用範囲に入るので、
\[
 \operatorname{Ext}^i(M,N)^*
 \simeq\operatorname{Ext}^{d-i}(N,M_\nu(-\ell)).
\]
$i>d$では右辺が負次数ゆえ零。
また$M=\pi A$、$N=M_\nu(-\ell)$と取ると
$\operatorname{Ext}^d(M,N)^*\simeq\operatorname{Hom}(N,N)\ne0$なので等号となる。
\end{proof}

\begin{lemma}[超越的中心値の有限大域次元]
\label{CSBGlem:global-transcendental}
$k_0\subset\CC$、$R_0$は$k_0$-代数、$z\in Z(R_0)$は$k_0$上超越的とする。
$R=R_0\otimes_{k_0}\CC$が左右noetherianで$\operatorname{gldim}R\le n$なら、
$k_0$上超越的な全ての$\lambda\in\CC$について
$\operatorname{gldim}R/(z-\lambda)\le n$である。
\end{lemma}
\begin{proof}
$\lambda$を含む$\CC/k_0$の超越基底$T$を選び、$K=k_0(T\setminus\{\lambda\})$とする。
体拡大は自由かつ忠実平坦で、$K$-線形な$\CC\to K$の$1\mapsto1$という分裂がある。
従って大域次元の上界とnoetherian性は$R_K=R_0\otimes K$へ降下する。
$z$の冪は$K$上でも独立なので、中心局所化
$R_\eta=R_K\otimes_{K[z]}K(z)$も大域次元高々$n$。
$z\mapsto\lambda$で$K(z)=K(\lambda)$とみなすと
$R_\eta\otimes_{K(z)}\CC=R/(z-\lambda)$である。
最後の体拡大は代数的分離拡大で、右辺は$R$の商ゆえnoetherian。
\cite[Lemmas~3.2, 3.3]{CSBGRSVBertini}の仮定が揃い、上界が保たれる。
任意の超越的体拡大が有限大域次元を保つと仮定したわけではない。
\end{proof}

\begin{theorem}[固定した良いambientでのvery-generalな正則性]
\label{CSBGthm:global-very-general}
$B$は中心一次元$u,v$を持つ左右noetherian AS正則五次元代数とし、
$G\in Z(B)_3$、$\operatorname{gr}_{x,y,z}B=S\otimes\CC[u,v]$とする。
$S$は整域、$\operatorname{gr}G=\gamma\in Z(S)_3$は非零正則元とする。
全ての$f\in\CC[u,v]_3$について$B/(G+f)$は$\CC[u,v]$上自由であり、
$\operatorname{qgr}(S/(\gamma))$は滑らかな楕円曲線のcoherent圏と仮定する。
有限表示の定義体$k_0$に対し、$a,d$が$k_0(b,c)$上代数的独立なら
\[
 \operatorname{gldim}\operatorname{qgr}
 B/(G+au^3+bu^2v+cuv^2+dv^3)=3.
\]
\end{theorem}
\begin{proof}
$D_u=B[u^{-1}]_0$はnoetherianで有限大域次元を持つ。
実際、中心局所化は大域次元を増やさず、
$B[u^{-1}]=D_u[u,u^{-1}]$から制限・誘導と評価$u=1$で上界を降下できる。
$t=v/u$とすると商の$u$-chartは
$D_u/(G/u^3+bt+ct^2+dt^3+a)$。
括弧内の$a$を除く中心元は$k_0(b,c,d)$上定義され、
最高主変数成分が非零正則なので、その冪は独立である。
前補題でchartの有限大域次元を得る。
$v$-chartも$d$を中心値として同様に扱える。
命題\ref{CSBGprop:global-two-charts}と
補題\ref{CSBGlem:gluing-serre-dimension}が結論を与える。
\end{proof}

\begin{lemma}[被約基底上の中心の持上げ]
\label{CSBGlem:global-center-lifting}
$H=B(E,\mathcal L,\tau)$は標数零の滑らかな楕円曲線のTHCRで、
$\mathcal L$はample、$\tau$は有限位数のtranslationとする。
$R$を被約有限型可換代数とし、filtered $R$-代数$U$が
$\operatorname{gr}U=H\otimes R$、各$F_dU$は有限自由で、
閉点への基底変更と両立するとする。
このとき各$z\in Z(H)_d$は$Z(U)\cap F_dU$へ持ち上がり、
$U$は有限型可換中心部分代数上有限となる。
\end{lemma}
\begin{proof}
斉次生成元の持上げ$u_i$を取り、有限自由加群間の交換子写像
$M_d:F_dU\to\bigoplus_iF_{d+\deg u_i}U$を作る。
Rains\cite[Theorem~2.13]{CSBGRainsElliptic}により各閉点$s$で
$\operatorname{gr}Z(U_s)=Z(H)$なので、$M_d$のrankは一定である。
被約性により高すぎるminorは恒等的に零となり、
最大rankのminorで被覆すれば$\ker M_d$は有限射影で、形成は基底変更と可換する。
最高次symbolの写像$\ker M_d\to Z(H)_d\otimes R$は
各閉点で全射、従ってNakayamaにより全射である。
右辺が自由なので、各斉次中心元は大域的に持ち上がる。
$H$は有限位数translationのnorm line bundleによる中心上有限である。
中心の有限生成元と中心上の有限加群生成元を持ち上げ、次数を下げる帰納法で
$U$も持ち上げ中心部分代数上有限と分かる。
ここではRainsの体上の定理を非被約基底へ拡張していない。
\end{proof}

\begin{theorem}[有限位数境界を持つ固定ambientの開な正則範囲]
\label{CSBGthm:global-open}
定理\ref{CSBGthm:global-very-general}の良い固定ambientを取り、さらに
$H=S/(\gamma)=B(E,\mathcal L,\tau)$のtranslationが有限位数とする。
ある非零多項式$P(a,b,c,d)$が存在して、$P\ne0$上では
中心三次商のtails大域次元が3となる。
周囲のPBW/Koszul解消とnoetherian AS条件が基底変更と両立する構成では、
この十分条件は任意の代数閉拡大体上でも成立する。
特に周囲の表示が有理数係数なら、有理数四係数の正則な商がZariski稠密に存在する。
\end{theorem}
\begin{proof}
$D_u=B[u^{-1}]_0$、$t=v/u$、$G_u=G/u^3$とする。
全係数を動かす$u$-chartの全代数は
\[
 \mathcal U=D_u[a,b,c,d]/(G_u+a+bt+ct^2+dt^3)
 \simeq D_u[b,c,d]
\]
なので左右noetherianかつ有限大域次元である。
$R=\CC[t,a,b,c,d]$をfiltration次数零とすると
$\operatorname{gr}\mathcal U=H\otimes R$で、各filtered pieceは有限自由。
前補題から$\mathcal U$は有限型可換中心部分代数$C_u$上有限である。
この$C_u$は係数環$R_0=\CC[a,b,c,d]$上も有限型である。
ここで\cite[Proposition~2.1, proof steps (6)--(7)]{CSBGRSVBertini}を
$(R,S,A)=(R_0,C_u,\mathcal U)$へ適用する。
その証明はある非零$P_u\in R_0$を逆にした上で、$\mathcal U$が$R_0$上平坦で、
相対対角が有限射影解消を持つことを与える。
従って全ての幾何学的fiberで有限大域次元が成立し、体拡大も許される。
$v$-chartでも$d$を消去して$P_v$を得る。
$P=P_uP_v$とし、二chartの貼合せとSerre双対性で結論する。
$\mathbb Q^4$は複素四次元affine空間でもZariski稠密なので、
非空開集合には有理数係数点が稠密に存在する。
\end{proof}

\begin{lemma}[位数6の具体的な楕円曲線計算]\label{CSBGlem:order-six}
$E_1=V(x^3+y^3+z^3+xyz)$で$O=[1:-1:0]$を原点とすると、
$p=[1:-1:-1]$は位数6であり、$2p=[1:0:-1]$である。
\end{lemma}
\begin{proof}
$E_1$は滑らかで、$O$は変曲点である。
$p$の接線$2x+y+z=0$へ代入すると
$7(y+z)(y-z)^2/8$となり、第三交点は$[0:1:-1]$である。
この点と$O,[1:0:-1]$は共線なので、接線の第三交点を反転すると表示した$2p$を得る。
$2p$は$O$と異なる変曲点なので位数3、$p$自体はHessianが非零で変曲点ではない。
従って$p$の位数は6である。各代入は現行の厳密検算に含まれる。
\end{proof}
$h=-1$の場合は$x,y$の交換が$S_{-1}$と$S_1$の次数付き同型を与えるので、境界translationの位数も同じである。

\begin{thebibliography}{99}
\bibitem{PairBW} S.~Barmeier and Z.~Wang,
\emph{Deformations of path algebras of quivers with relations},
\url{https://arxiv.org/abs/2002.10001}, v4 (2021), Theorem 10.3 and Lemma 10.9.
\bibitem{RainsMagma} E.~Rains, central-element computations in MAGMA form,
\url{https://math.mit.edu/~etingof/delpezzocenter}; E6 data.
\bibitem{AS}
M.~Artin and W.~F.~Schelter,
\textit{Graded algebras of global dimension 3},
Adv. Math. \textbf{66} (1987), 171--216.

\bibitem{AZ}
M.~Artin and J.~J.~Zhang,
\textit{Noncommutative projective schemes},
Adv. Math. \textbf{109} (1994), 228--287.
\href{https://doi.org/10.1006/aima.1994.1087}{doi:10.1006/aima.1994.1087}.

\bibitem{ATV}
M.~Artin, J.~Tate and M.~Van den Bergh,
\textit{Some algebras associated to automorphisms of elliptic curves},
in \textit{The Grothendieck Festschrift}, Vol.~I,
Birkh\"auser (1990), 33--85.
\href{https://doi.org/10.1007/978-0-8176-4574-8_3}{doi:10.1007/978-0-8176-4574-8\_3}.

\bibitem{AD}
M.~Artebani and I.~Dolgachev,
\textit{The Hesse pencil of plane cubic curves},
Enseign. Math. \textbf{55} (2009), 235--273,
\href{https://doi.org/10.4171/LEM/55-3-3}{doi:10.4171/LEM/55-3-3}.

\bibitem{BT}
R.~Berger and R.~Taillefer,
\textit{Poincar\'e--Birkhoff--Witt deformations of Calabi--Yau algebras},
J. Noncommut. Geom. \textbf{1} (2007), 241--270,
\href{https://doi.org/10.4171/JNCG/6}{doi:10.4171/JNCG/6},
\href{https://arxiv.org/abs/math/0610112}{arXiv:math/0610112}.

\bibitem{CKS}
A.~Chirvasitu, R.~Kanda and S.~P.~Smith,
\textit{New Artin--Schelter regular and Calabi--Yau algebras via normal extensions},
\href{https://arxiv.org/abs/1706.05754}{arXiv:1706.05754v3},
\href{https://doi.org/10.1090/tran/7672}{doi:10.1090/tran/7672}.

\bibitem{EG}
P.~Etingof and V.~Ginzburg,
\textit{Noncommutative del Pezzo surfaces and Calabi--Yau algebras},
J. Eur. Math. Soc. \textbf{12} (2010), 1371--1416,
\href{https://doi.org/10.4171/JEMS/235}{doi:10.4171/JEMS/235},
\href{https://arxiv.org/abs/0709.3593}{arXiv:0709.3593v4}.

\bibitem{LPT}
F.~Loray, J.~V.~Pereira and F.~Touzet,
\textit{Foliations with trivial canonical bundle on Fano 3-folds},
Math. Nachr. \textbf{286} (2013), 921--940,
\href{https://arxiv.org/abs/1112.3871}{arXiv:1112.3871}.

\bibitem{Mayanskiy}
E.~Mayanskiy,
\textit{Poisson cohomology of two Fano threefolds},
J. Algebra \textbf{424} (2015), 21--45,
\href{https://doi.org/10.1016/j.jalgebra.2014.08.049}{doi:10.1016/j.jalgebra.2014.08.049},
\href{https://arxiv.org/abs/1303.5916}{arXiv:1303.5916}.

\bibitem{MS}
I.~Mori and S.~P.~Smith,
\textit{The classification of 3-Calabi--Yau algebras with 3 generators and 3 quadratic relations},
Math. Z. (2017),
\href{https://doi.org/10.1007/s00209-016-1824-5}{doi:10.1007/s00209-016-1824-5},
\href{https://arxiv.org/abs/1502.07403}{arXiv:1502.07403v2}.

\bibitem{Pym}
B.~Pym,
\textit{Quantum deformations of projective three-space},
Adv. Math. \textbf{281} (2015), 1216--1241,
\href{https://doi.org/10.1016/j.aim.2015.06.005}{doi:10.1016/j.aim.2015.06.005},
\href{https://arxiv.org/abs/1403.6444}{arXiv:1403.6444}.

\bibitem{Rains}
E.~M.~Rains,
\textit{Filtered deformations of elliptic algebras},
\href{https://arxiv.org/abs/2107.13540}{arXiv:2107.13540v2},
\S2, Theorem~2.13.

\bibitem{RSVdB}
J.~V.~Rennemo, E.~Segal and M.~Van den Bergh,
\textit{A non-commutative Bertini theorem},
J. Noncommut. Geom. \textbf{13} (2019), 609--616,
\href{https://doi.org/10.4171/JNCG/334}{doi:10.4171/JNCG/334},
\href{https://arxiv.org/abs/1705.01366}{arXiv:1705.01366v2}.

\bibitem{Rosenlicht}
M.~Rosenlicht,
\textit{Some basic theorems on algebraic groups},
Amer. J. Math. \textbf{78} (1956), 401--443.

\bibitem{YZ}
A.~Yekutieli and J.~J.~Zhang,
\textit{Serre duality for noncommutative projective schemes},
Proc. Amer. Math. Soc. \textbf{125} (1997), 697--707,
\href{https://doi.org/10.1090/S0002-9939-97-03782-9}{doi:10.1090/S0002-9939-97-03782-9}.

\bibitem{ZSL}
G.-S.~Zhou, Y.~Shen and D.-M.~Lu,
\textit{Skew Calabi--Yau property of normal extensions},
Manuscripta Math. \textbf{161} (2020), 125--140,
\href{https://doi.org/10.1007/s00229-018-1064-6}{doi:10.1007/s00229-018-1064-6},
\href{https://arxiv.org/abs/1709.09292}{arXiv:1709.09292}.

\bibitem{Zhang}
J.~J.~Zhang,
\textit{Twisted graded algebras and equivalences of graded categories},
Proc. London Math. Soc. \textbf{72} (1996), 281--311,
\href{https://doi.org/10.1112/plms/s3-72.2.281}{doi:10.1112/plms/s3-72.2.281}.
\bibitem{PBondalPolishchukHelices}
A.~I.~Bondal and A.~E.~Polishchuk,
\emph{Homological properties of associative algebras: the method of helices},
Russian Acad. Sci. Izv. Math. \textbf{42} (1994), no.~2, 219--260.
Proposition~4.1 and \S5 (Koszul--Frobenius duality and periodicity).
\href{https://doi.org/10.1070/IM1994v042n02ABEH001536}
{doi:10.1070/IM1994v042n02ABEH001536}.

\bibitem{PChanNymanNormalFamily}
D.~Chan and A.~Nyman,
\emph{Noncommutative linear systems and noncommutative elliptic curves},
\href{https://arxiv.org/abs/2211.12465}{arXiv:2211.12465} (2022).
Definition~8.6 (normal families of elements in indexed algebras).

\bibitem{CSBGCCKM}
T.~Cassidy, A.~Conner, E.~Kirkman and W.~F.~Moore,
\emph{Periodic free resolutions from twisted matrix factorizations},
J. Algebra 455 (2016), 137--163;
\url{https://arxiv.org/abs/1307.8415}, Theorems~4.7 and 5.5.

\bibitem{CSBGRSVBertini}
J.~V.~Rennemo, E.~Segal and M.~Van den Bergh,
\emph{A non-commutative Bertini theorem},
\url{https://arxiv.org/abs/1705.01366}, v2 (2020), Lemmas~3.4 and 3.5.

\bibitem{CSBGChenGeneralizedSerre}
X.-W.~Chen, \emph{Generalized Serre duality},
\url{https://arxiv.org/abs/math/0610258}, Theorem~3.10.

\bibitem{CSBGRainsElliptic}
E.~M.~Rains, \emph{Filtered deformations of elliptic algebras},
\url{https://arxiv.org/abs/2107.13540}, v2 (2022), Theorem~2.13 and its preceding discussion.

\bibitem{CSBGYZSerre}
A. Yekutieli and J. J. Zhang,
\emph{Serre duality for noncommutative projective schemes},
Proc. Amer. Math. Soc. 125 (1997), 697--707;
\url{https://doi.org/10.1090/S0002-9939-97-03782-9}.

\end{thebibliography}
\end{document}